% Options for packages loaded elsewhere
% Options for packages loaded elsewhere
\PassOptionsToPackage{unicode}{hyperref}
\PassOptionsToPackage{hyphens}{url}
\PassOptionsToPackage{dvipsnames,svgnames,x11names}{xcolor}
%
\documentclass[
  letterpaper,
  DIV=11,
  numbers=noendperiod]{scrartcl}
\usepackage{xcolor}
\usepackage{amsmath,amssymb}
\setcounter{secnumdepth}{-\maxdimen} % remove section numbering
\usepackage{iftex}
\ifPDFTeX
  \usepackage[T1]{fontenc}
  \usepackage[utf8]{inputenc}
  \usepackage{textcomp} % provide euro and other symbols
\else % if luatex or xetex
  \usepackage{unicode-math} % this also loads fontspec
  \defaultfontfeatures{Scale=MatchLowercase}
  \defaultfontfeatures[\rmfamily]{Ligatures=TeX,Scale=1}
\fi
\usepackage{lmodern}
\ifPDFTeX\else
  % xetex/luatex font selection
\fi
% Use upquote if available, for straight quotes in verbatim environments
\IfFileExists{upquote.sty}{\usepackage{upquote}}{}
\IfFileExists{microtype.sty}{% use microtype if available
  \usepackage[]{microtype}
  \UseMicrotypeSet[protrusion]{basicmath} % disable protrusion for tt fonts
}{}
\makeatletter
\@ifundefined{KOMAClassName}{% if non-KOMA class
  \IfFileExists{parskip.sty}{%
    \usepackage{parskip}
  }{% else
    \setlength{\parindent}{0pt}
    \setlength{\parskip}{6pt plus 2pt minus 1pt}}
}{% if KOMA class
  \KOMAoptions{parskip=half}}
\makeatother
% Make \paragraph and \subparagraph free-standing
\makeatletter
\ifx\paragraph\undefined\else
  \let\oldparagraph\paragraph
  \renewcommand{\paragraph}{
    \@ifstar
      \xxxParagraphStar
      \xxxParagraphNoStar
  }
  \newcommand{\xxxParagraphStar}[1]{\oldparagraph*{#1}\mbox{}}
  \newcommand{\xxxParagraphNoStar}[1]{\oldparagraph{#1}\mbox{}}
\fi
\ifx\subparagraph\undefined\else
  \let\oldsubparagraph\subparagraph
  \renewcommand{\subparagraph}{
    \@ifstar
      \xxxSubParagraphStar
      \xxxSubParagraphNoStar
  }
  \newcommand{\xxxSubParagraphStar}[1]{\oldsubparagraph*{#1}\mbox{}}
  \newcommand{\xxxSubParagraphNoStar}[1]{\oldsubparagraph{#1}\mbox{}}
\fi
\makeatother

\usepackage{color}
\usepackage{fancyvrb}
\newcommand{\VerbBar}{|}
\newcommand{\VERB}{\Verb[commandchars=\\\{\}]}
\DefineVerbatimEnvironment{Highlighting}{Verbatim}{commandchars=\\\{\}}
% Add ',fontsize=\small' for more characters per line
\usepackage{framed}
\definecolor{shadecolor}{RGB}{241,243,245}
\newenvironment{Shaded}{\begin{snugshade}}{\end{snugshade}}
\newcommand{\AlertTok}[1]{\textcolor[rgb]{0.68,0.00,0.00}{#1}}
\newcommand{\AnnotationTok}[1]{\textcolor[rgb]{0.37,0.37,0.37}{#1}}
\newcommand{\AttributeTok}[1]{\textcolor[rgb]{0.40,0.45,0.13}{#1}}
\newcommand{\BaseNTok}[1]{\textcolor[rgb]{0.68,0.00,0.00}{#1}}
\newcommand{\BuiltInTok}[1]{\textcolor[rgb]{0.00,0.23,0.31}{#1}}
\newcommand{\CharTok}[1]{\textcolor[rgb]{0.13,0.47,0.30}{#1}}
\newcommand{\CommentTok}[1]{\textcolor[rgb]{0.37,0.37,0.37}{#1}}
\newcommand{\CommentVarTok}[1]{\textcolor[rgb]{0.37,0.37,0.37}{\textit{#1}}}
\newcommand{\ConstantTok}[1]{\textcolor[rgb]{0.56,0.35,0.01}{#1}}
\newcommand{\ControlFlowTok}[1]{\textcolor[rgb]{0.00,0.23,0.31}{\textbf{#1}}}
\newcommand{\DataTypeTok}[1]{\textcolor[rgb]{0.68,0.00,0.00}{#1}}
\newcommand{\DecValTok}[1]{\textcolor[rgb]{0.68,0.00,0.00}{#1}}
\newcommand{\DocumentationTok}[1]{\textcolor[rgb]{0.37,0.37,0.37}{\textit{#1}}}
\newcommand{\ErrorTok}[1]{\textcolor[rgb]{0.68,0.00,0.00}{#1}}
\newcommand{\ExtensionTok}[1]{\textcolor[rgb]{0.00,0.23,0.31}{#1}}
\newcommand{\FloatTok}[1]{\textcolor[rgb]{0.68,0.00,0.00}{#1}}
\newcommand{\FunctionTok}[1]{\textcolor[rgb]{0.28,0.35,0.67}{#1}}
\newcommand{\ImportTok}[1]{\textcolor[rgb]{0.00,0.46,0.62}{#1}}
\newcommand{\InformationTok}[1]{\textcolor[rgb]{0.37,0.37,0.37}{#1}}
\newcommand{\KeywordTok}[1]{\textcolor[rgb]{0.00,0.23,0.31}{\textbf{#1}}}
\newcommand{\NormalTok}[1]{\textcolor[rgb]{0.00,0.23,0.31}{#1}}
\newcommand{\OperatorTok}[1]{\textcolor[rgb]{0.37,0.37,0.37}{#1}}
\newcommand{\OtherTok}[1]{\textcolor[rgb]{0.00,0.23,0.31}{#1}}
\newcommand{\PreprocessorTok}[1]{\textcolor[rgb]{0.68,0.00,0.00}{#1}}
\newcommand{\RegionMarkerTok}[1]{\textcolor[rgb]{0.00,0.23,0.31}{#1}}
\newcommand{\SpecialCharTok}[1]{\textcolor[rgb]{0.37,0.37,0.37}{#1}}
\newcommand{\SpecialStringTok}[1]{\textcolor[rgb]{0.13,0.47,0.30}{#1}}
\newcommand{\StringTok}[1]{\textcolor[rgb]{0.13,0.47,0.30}{#1}}
\newcommand{\VariableTok}[1]{\textcolor[rgb]{0.07,0.07,0.07}{#1}}
\newcommand{\VerbatimStringTok}[1]{\textcolor[rgb]{0.13,0.47,0.30}{#1}}
\newcommand{\WarningTok}[1]{\textcolor[rgb]{0.37,0.37,0.37}{\textit{#1}}}

\usepackage{longtable,booktabs,array}
\usepackage{calc} % for calculating minipage widths
% Correct order of tables after \paragraph or \subparagraph
\usepackage{etoolbox}
\makeatletter
\patchcmd\longtable{\par}{\if@noskipsec\mbox{}\fi\par}{}{}
\makeatother
% Allow footnotes in longtable head/foot
\IfFileExists{footnotehyper.sty}{\usepackage{footnotehyper}}{\usepackage{footnote}}
\makesavenoteenv{longtable}
\usepackage{graphicx}
\makeatletter
\newsavebox\pandoc@box
\newcommand*\pandocbounded[1]{% scales image to fit in text height/width
  \sbox\pandoc@box{#1}%
  \Gscale@div\@tempa{\textheight}{\dimexpr\ht\pandoc@box+\dp\pandoc@box\relax}%
  \Gscale@div\@tempb{\linewidth}{\wd\pandoc@box}%
  \ifdim\@tempb\p@<\@tempa\p@\let\@tempa\@tempb\fi% select the smaller of both
  \ifdim\@tempa\p@<\p@\scalebox{\@tempa}{\usebox\pandoc@box}%
  \else\usebox{\pandoc@box}%
  \fi%
}
% Set default figure placement to htbp
\def\fps@figure{htbp}
\makeatother





\setlength{\emergencystretch}{3em} % prevent overfull lines

\providecommand{\tightlist}{%
  \setlength{\itemsep}{0pt}\setlength{\parskip}{0pt}}



 


\KOMAoption{captions}{tableheading}
\makeatletter
\@ifpackageloaded{caption}{}{\usepackage{caption}}
\AtBeginDocument{%
\ifdefined\contentsname
  \renewcommand*\contentsname{Table of contents}
\else
  \newcommand\contentsname{Table of contents}
\fi
\ifdefined\listfigurename
  \renewcommand*\listfigurename{List of Figures}
\else
  \newcommand\listfigurename{List of Figures}
\fi
\ifdefined\listtablename
  \renewcommand*\listtablename{List of Tables}
\else
  \newcommand\listtablename{List of Tables}
\fi
\ifdefined\figurename
  \renewcommand*\figurename{Figure}
\else
  \newcommand\figurename{Figure}
\fi
\ifdefined\tablename
  \renewcommand*\tablename{Table}
\else
  \newcommand\tablename{Table}
\fi
}
\@ifpackageloaded{float}{}{\usepackage{float}}
\floatstyle{ruled}
\@ifundefined{c@chapter}{\newfloat{codelisting}{h}{lop}}{\newfloat{codelisting}{h}{lop}[chapter]}
\floatname{codelisting}{Listing}
\newcommand*\listoflistings{\listof{codelisting}{List of Listings}}
\makeatother
\makeatletter
\makeatother
\makeatletter
\@ifpackageloaded{caption}{}{\usepackage{caption}}
\@ifpackageloaded{subcaption}{}{\usepackage{subcaption}}
\makeatother
\usepackage{bookmark}
\IfFileExists{xurl.sty}{\usepackage{xurl}}{} % add URL line breaks if available
\urlstyle{same}
\hypersetup{
  pdftitle={BASIC COMMANDS (Assignment-3)},
  pdfauthor={Tushar Sharma},
  colorlinks=true,
  linkcolor={blue},
  filecolor={Maroon},
  citecolor={Blue},
  urlcolor={Blue},
  pdfcreator={LaTeX via pandoc}}


\title{BASIC COMMANDS (Assignment-3)}
\author{Tushar Sharma}
\date{}
\begin{document}
\maketitle


\subsection{Practicing Basic Operations of
R}\label{practicing-basic-operations-of-r}

\subsubsection{Create a vector - c()}\label{create-a-vector---c}

\begin{Shaded}
\begin{Highlighting}[]
\CommentTok{\#Vectors of number }
\NormalTok{v1 }\OtherTok{\textless{}{-}} \FunctionTok{c}\NormalTok{(}\DecValTok{1}\NormalTok{,}\DecValTok{2}\NormalTok{,}\DecValTok{3}\NormalTok{,}\DecValTok{4}\NormalTok{,}\DecValTok{5}\NormalTok{,}\DecValTok{6}\NormalTok{)}
\FunctionTok{print}\NormalTok{(v1)}
\end{Highlighting}
\end{Shaded}

\begin{verbatim}
[1] 1 2 3 4 5 6
\end{verbatim}

\begin{Shaded}
\begin{Highlighting}[]
\CommentTok{\#vectorss of characters}
\NormalTok{v2 }\OtherTok{\textless{}{-}} \FunctionTok{c}\NormalTok{(}\StringTok{"Anuj "}\NormalTok{, }\StringTok{"Tushar"}\NormalTok{, }\StringTok{"datta "}\NormalTok{)}
\FunctionTok{print}\NormalTok{(v2)}
\end{Highlighting}
\end{Shaded}

\begin{verbatim}
[1] "Anuj "  "Tushar" "datta "
\end{verbatim}

\begin{Shaded}
\begin{Highlighting}[]
\CommentTok{\#missalaneous vector}
\NormalTok{v3 }\OtherTok{\textless{}{-}} \FunctionTok{c}\NormalTok{(}\DecValTok{1}\NormalTok{,}\StringTok{"hello"}\NormalTok{, }\ConstantTok{FALSE}\NormalTok{)}
\FunctionTok{print}\NormalTok{(v3)}
\end{Highlighting}
\end{Shaded}

\begin{verbatim}
[1] "1"     "hello" "FALSE"
\end{verbatim}

\subsubsection{Generating sequences of number -
seq()}\label{generating-sequences-of-number---seq}

\begin{Shaded}
\begin{Highlighting}[]
\NormalTok{s1 }\OtherTok{\textless{}{-}} \FunctionTok{seq}\NormalTok{(}\DecValTok{1}\NormalTok{, }\DecValTok{10}\NormalTok{)}
\FunctionTok{print}\NormalTok{(s1)}
\end{Highlighting}
\end{Shaded}

\begin{verbatim}
 [1]  1  2  3  4  5  6  7  8  9 10
\end{verbatim}

\begin{Shaded}
\begin{Highlighting}[]
\NormalTok{s2 }\OtherTok{\textless{}{-}} \FunctionTok{seq}\NormalTok{(}\DecValTok{1}\NormalTok{, }\DecValTok{10}\NormalTok{, }\AttributeTok{by =} \DecValTok{2}\NormalTok{)}
\FunctionTok{print}\NormalTok{(s2)}
\end{Highlighting}
\end{Shaded}

\begin{verbatim}
[1] 1 3 5 7 9
\end{verbatim}

\begin{Shaded}
\begin{Highlighting}[]
\CommentTok{\#equally spaced}
\NormalTok{s3 }\OtherTok{\textless{}{-}} \FunctionTok{seq}\NormalTok{(}\DecValTok{1}\NormalTok{, }\DecValTok{10}\NormalTok{, }\AttributeTok{length.out =} \DecValTok{5}\NormalTok{)}
\FunctionTok{print}\NormalTok{(s3)}
\end{Highlighting}
\end{Shaded}

\begin{verbatim}
[1]  1.00  3.25  5.50  7.75 10.00
\end{verbatim}

\begin{Shaded}
\begin{Highlighting}[]
\DocumentationTok{\#\#\# Length of vector {-} length()}

\NormalTok{l1 }\OtherTok{\textless{}{-}} \FunctionTok{length}\NormalTok{(v1)}
\FunctionTok{print}\NormalTok{(l1)}
\end{Highlighting}
\end{Shaded}

\begin{verbatim}
[1] 6
\end{verbatim}

\begin{Shaded}
\begin{Highlighting}[]
\NormalTok{l2 }\OtherTok{\textless{}{-}} \FunctionTok{length}\NormalTok{(v2)}
\FunctionTok{print}\NormalTok{(l2)}
\end{Highlighting}
\end{Shaded}

\begin{verbatim}
[1] 3
\end{verbatim}

\begin{Shaded}
\begin{Highlighting}[]
\NormalTok{l3 }\OtherTok{\textless{}{-}} \FunctionTok{length}\NormalTok{(v3)}
\FunctionTok{print}\NormalTok{(l3)}
\end{Highlighting}
\end{Shaded}

\begin{verbatim}
[1] 3
\end{verbatim}

\subsubsection{Sorting values - sort()}\label{sorting-values---sort}

\begin{Shaded}
\begin{Highlighting}[]
\NormalTok{so1 }\OtherTok{\textless{}{-}} \FunctionTok{sort}\NormalTok{(}\FunctionTok{c}\NormalTok{(}\DecValTok{5}\NormalTok{, }\DecValTok{2}\NormalTok{, }\DecValTok{9}\NormalTok{, }\DecValTok{1}\NormalTok{))}
\FunctionTok{print}\NormalTok{(so1)}
\end{Highlighting}
\end{Shaded}

\begin{verbatim}
[1] 1 2 5 9
\end{verbatim}

\begin{Shaded}
\begin{Highlighting}[]
\NormalTok{so2 }\OtherTok{\textless{}{-}} \FunctionTok{sort}\NormalTok{(}\FunctionTok{c}\NormalTok{(}\DecValTok{10}\NormalTok{, }\DecValTok{4}\NormalTok{, }\DecValTok{7}\NormalTok{), }\AttributeTok{decreasing =} \ConstantTok{TRUE}\NormalTok{)}
\FunctionTok{print}\NormalTok{(so2)}
\end{Highlighting}
\end{Shaded}

\begin{verbatim}
[1] 10  7  4
\end{verbatim}

\begin{Shaded}
\begin{Highlighting}[]
\NormalTok{so3 }\OtherTok{\textless{}{-}} \FunctionTok{sort}\NormalTok{(}\FunctionTok{c}\NormalTok{(}\DecValTok{3}\NormalTok{, }\DecValTok{3}\NormalTok{, }\DecValTok{1}\NormalTok{, }\DecValTok{2}\NormalTok{))}
\FunctionTok{print}\NormalTok{(so3)}
\end{Highlighting}
\end{Shaded}

\begin{verbatim}
[1] 1 2 3 3
\end{verbatim}

\subsubsection{Ordering elements -
order()}\label{ordering-elements---order}

\begin{Shaded}
\begin{Highlighting}[]
\NormalTok{o1 }\OtherTok{\textless{}{-}} \FunctionTok{order}\NormalTok{(}\FunctionTok{c}\NormalTok{(}\DecValTok{40}\NormalTok{, }\DecValTok{10}\NormalTok{, }\DecValTok{30}\NormalTok{))}
\FunctionTok{print}\NormalTok{(o1)}
\end{Highlighting}
\end{Shaded}

\begin{verbatim}
[1] 2 3 1
\end{verbatim}

\begin{Shaded}
\begin{Highlighting}[]
\NormalTok{o2 }\OtherTok{\textless{}{-}}\NormalTok{ v1[}\FunctionTok{order}\NormalTok{(v1)]}
\FunctionTok{print}\NormalTok{(o2)}
\end{Highlighting}
\end{Shaded}

\begin{verbatim}
[1] 1 2 3 4 5 6
\end{verbatim}

\begin{Shaded}
\begin{Highlighting}[]
\NormalTok{o3 }\OtherTok{\textless{}{-}}\NormalTok{ v1[}\FunctionTok{order}\NormalTok{(v1, }\AttributeTok{decreasing =} \ConstantTok{TRUE}\NormalTok{)]}
\FunctionTok{print}\NormalTok{(o3)}
\end{Highlighting}
\end{Shaded}

\begin{verbatim}
[1] 6 5 4 3 2 1
\end{verbatim}

\subsubsection{Reverse a vector - rev()}\label{reverse-a-vector---rev}

\begin{Shaded}
\begin{Highlighting}[]
\NormalTok{r1 }\OtherTok{\textless{}{-}} \FunctionTok{rev}\NormalTok{(v1)}
\FunctionTok{print}\NormalTok{(r1)}
\end{Highlighting}
\end{Shaded}

\begin{verbatim}
[1] 6 5 4 3 2 1
\end{verbatim}

\begin{Shaded}
\begin{Highlighting}[]
\NormalTok{r2 }\OtherTok{\textless{}{-}} \FunctionTok{rev}\NormalTok{(v2)}
\FunctionTok{print}\NormalTok{(r2)}
\end{Highlighting}
\end{Shaded}

\begin{verbatim}
[1] "datta " "Tushar" "Anuj " 
\end{verbatim}

\begin{Shaded}
\begin{Highlighting}[]
\NormalTok{r3 }\OtherTok{\textless{}{-}} \FunctionTok{rev}\NormalTok{(}\FunctionTok{c}\NormalTok{(}\ConstantTok{TRUE}\NormalTok{, }\ConstantTok{FALSE}\NormalTok{, }\ConstantTok{TRUE}\NormalTok{))}
\FunctionTok{print}\NormalTok{(r3)}
\end{Highlighting}
\end{Shaded}

\begin{verbatim}
[1]  TRUE FALSE  TRUE
\end{verbatim}

\subsubsection{Unique values - unique()}\label{unique-values---unique}

\begin{Shaded}
\begin{Highlighting}[]
\NormalTok{u1 }\OtherTok{\textless{}{-}} \FunctionTok{unique}\NormalTok{(}\FunctionTok{c}\NormalTok{(}\DecValTok{1}\NormalTok{, }\DecValTok{1}\NormalTok{, }\DecValTok{2}\NormalTok{, }\DecValTok{3}\NormalTok{))}
\FunctionTok{print}\NormalTok{(u1)}
\end{Highlighting}
\end{Shaded}

\begin{verbatim}
[1] 1 2 3
\end{verbatim}

\begin{Shaded}
\begin{Highlighting}[]
\NormalTok{u2 }\OtherTok{\textless{}{-}} \FunctionTok{unique}\NormalTok{(}\FunctionTok{c}\NormalTok{(}\StringTok{"a"}\NormalTok{, }\StringTok{"b"}\NormalTok{, }\StringTok{"a"}\NormalTok{))}
\FunctionTok{print}\NormalTok{(u2)}
\end{Highlighting}
\end{Shaded}

\begin{verbatim}
[1] "a" "b"
\end{verbatim}

\begin{Shaded}
\begin{Highlighting}[]
\NormalTok{u3 }\OtherTok{\textless{}{-}} \FunctionTok{unique}\NormalTok{(}\FunctionTok{c}\NormalTok{(}\ConstantTok{TRUE}\NormalTok{, }\ConstantTok{TRUE}\NormalTok{, }\ConstantTok{FALSE}\NormalTok{))}
\FunctionTok{print}\NormalTok{(u3)}
\end{Highlighting}
\end{Shaded}

\begin{verbatim}
[1]  TRUE FALSE
\end{verbatim}

\subsubsection{Frequency table - table()}\label{frequency-table---table}

\begin{Shaded}
\begin{Highlighting}[]
\NormalTok{t1 }\OtherTok{\textless{}{-}} \FunctionTok{table}\NormalTok{(}\FunctionTok{c}\NormalTok{(}\DecValTok{1}\NormalTok{, }\DecValTok{1}\NormalTok{, }\DecValTok{2}\NormalTok{, }\DecValTok{3}\NormalTok{))}
\FunctionTok{print}\NormalTok{(t1)}
\end{Highlighting}
\end{Shaded}

\begin{verbatim}

1 2 3 
2 1 1 
\end{verbatim}

\begin{Shaded}
\begin{Highlighting}[]
\NormalTok{t2 }\OtherTok{\textless{}{-}} \FunctionTok{table}\NormalTok{(v2)}
\FunctionTok{print}\NormalTok{(t2)}
\end{Highlighting}
\end{Shaded}

\begin{verbatim}
v2
 Anuj  datta  Tushar 
     1      1      1 
\end{verbatim}

\begin{Shaded}
\begin{Highlighting}[]
\NormalTok{t3 }\OtherTok{\textless{}{-}} \FunctionTok{table}\NormalTok{(}\FunctionTok{c}\NormalTok{(}\StringTok{"yes"}\NormalTok{, }\StringTok{"no"}\NormalTok{, }\StringTok{"yes"}\NormalTok{))}
\FunctionTok{print}\NormalTok{(t3)}
\end{Highlighting}
\end{Shaded}

\begin{verbatim}

 no yes 
  1   2 
\end{verbatim}

\subsubsection{Match values - match()}\label{match-values---match}

\begin{Shaded}
\begin{Highlighting}[]
\NormalTok{m1 }\OtherTok{\textless{}{-}} \FunctionTok{match}\NormalTok{(}\DecValTok{2}\NormalTok{, v1)}
\FunctionTok{print}\NormalTok{(m1)}
\end{Highlighting}
\end{Shaded}

\begin{verbatim}
[1] 2
\end{verbatim}

\begin{Shaded}
\begin{Highlighting}[]
\NormalTok{m2 }\OtherTok{\textless{}{-}} \FunctionTok{match}\NormalTok{(}\StringTok{"Tushar"}\NormalTok{, v2)}
\FunctionTok{print}\NormalTok{(m2)}
\end{Highlighting}
\end{Shaded}

\begin{verbatim}
[1] 2
\end{verbatim}

\begin{Shaded}
\begin{Highlighting}[]
\NormalTok{m3 }\OtherTok{\textless{}{-}} \FunctionTok{match}\NormalTok{(}\DecValTok{10}\NormalTok{, v1)}
\FunctionTok{print}\NormalTok{(m3)}
\end{Highlighting}
\end{Shaded}

\begin{verbatim}
[1] NA
\end{verbatim}

\subsubsection{Membership test - \%in\%}\label{membership-test---in}

\begin{Shaded}
\begin{Highlighting}[]
\NormalTok{in1 }\OtherTok{\textless{}{-}} \DecValTok{2} \SpecialCharTok{\%in\%}\NormalTok{ v1}
\FunctionTok{print}\NormalTok{(in1)}
\end{Highlighting}
\end{Shaded}

\begin{verbatim}
[1] TRUE
\end{verbatim}

\begin{Shaded}
\begin{Highlighting}[]
\NormalTok{in2 }\OtherTok{\textless{}{-}} \StringTok{"Anuj "} \SpecialCharTok{\%in\%}\NormalTok{ v2}
\FunctionTok{print}\NormalTok{(in2)}
\end{Highlighting}
\end{Shaded}

\begin{verbatim}
[1] TRUE
\end{verbatim}

\begin{Shaded}
\begin{Highlighting}[]
\NormalTok{in3 }\OtherTok{\textless{}{-}} \DecValTok{100} \SpecialCharTok{\%in\%}\NormalTok{ v1}
\FunctionTok{print}\NormalTok{(in3)}
\end{Highlighting}
\end{Shaded}

\begin{verbatim}
[1] FALSE
\end{verbatim}

\subsubsection{Find positions - which()}\label{find-positions---which}

\begin{Shaded}
\begin{Highlighting}[]
\NormalTok{w1 }\OtherTok{\textless{}{-}} \FunctionTok{which}\NormalTok{(v1 }\SpecialCharTok{\textgreater{}} \DecValTok{3}\NormalTok{)}
\FunctionTok{print}\NormalTok{(w1)}
\end{Highlighting}
\end{Shaded}

\begin{verbatim}
[1] 4 5 6
\end{verbatim}

\begin{Shaded}
\begin{Highlighting}[]
\NormalTok{w2 }\OtherTok{\textless{}{-}} \FunctionTok{which}\NormalTok{(v1 }\SpecialCharTok{==} \DecValTok{5}\NormalTok{)}
\FunctionTok{print}\NormalTok{(w2)}
\end{Highlighting}
\end{Shaded}

\begin{verbatim}
[1] 5
\end{verbatim}

\begin{Shaded}
\begin{Highlighting}[]
\NormalTok{w3 }\OtherTok{\textless{}{-}} \FunctionTok{which}\NormalTok{(}\FunctionTok{c}\NormalTok{(}\ConstantTok{TRUE}\NormalTok{, }\ConstantTok{FALSE}\NormalTok{, }\ConstantTok{TRUE}\NormalTok{))}
\FunctionTok{print}\NormalTok{(w3)}
\end{Highlighting}
\end{Shaded}

\begin{verbatim}
[1] 1 3
\end{verbatim}

\subsubsection{Any condition TRUE -
any()}\label{any-condition-true---any}

\begin{Shaded}
\begin{Highlighting}[]
\NormalTok{a1 }\OtherTok{\textless{}{-}} \FunctionTok{any}\NormalTok{(v1 }\SpecialCharTok{\textgreater{}} \DecValTok{5}\NormalTok{)}
\FunctionTok{print}\NormalTok{(a1)}
\end{Highlighting}
\end{Shaded}

\begin{verbatim}
[1] TRUE
\end{verbatim}

\begin{Shaded}
\begin{Highlighting}[]
\NormalTok{a2 }\OtherTok{\textless{}{-}} \FunctionTok{any}\NormalTok{(v1 }\SpecialCharTok{\textless{}} \DecValTok{0}\NormalTok{)}
\FunctionTok{print}\NormalTok{(a2)}
\end{Highlighting}
\end{Shaded}

\begin{verbatim}
[1] FALSE
\end{verbatim}

\begin{Shaded}
\begin{Highlighting}[]
\NormalTok{a3 }\OtherTok{\textless{}{-}} \FunctionTok{any}\NormalTok{(}\FunctionTok{is.na}\NormalTok{(}\FunctionTok{c}\NormalTok{(}\DecValTok{1}\NormalTok{, }\ConstantTok{NA}\NormalTok{, }\DecValTok{3}\NormalTok{)))}
\FunctionTok{print}\NormalTok{(a3)}
\end{Highlighting}
\end{Shaded}

\begin{verbatim}
[1] TRUE
\end{verbatim}

\subsubsection{All conditions TRUE -
all()}\label{all-conditions-true---all}

\begin{Shaded}
\begin{Highlighting}[]
\NormalTok{al1 }\OtherTok{\textless{}{-}} \FunctionTok{all}\NormalTok{(v1 }\SpecialCharTok{\textgreater{}} \DecValTok{0}\NormalTok{)}
\FunctionTok{print}\NormalTok{(al1)}
\end{Highlighting}
\end{Shaded}

\begin{verbatim}
[1] TRUE
\end{verbatim}

\begin{Shaded}
\begin{Highlighting}[]
\NormalTok{al2 }\OtherTok{\textless{}{-}} \FunctionTok{all}\NormalTok{(v1 }\SpecialCharTok{\textgreater{}} \DecValTok{3}\NormalTok{)}
\FunctionTok{print}\NormalTok{(al2)}
\end{Highlighting}
\end{Shaded}

\begin{verbatim}
[1] FALSE
\end{verbatim}

\begin{Shaded}
\begin{Highlighting}[]
\NormalTok{al3 }\OtherTok{\textless{}{-}} \FunctionTok{all}\NormalTok{(}\FunctionTok{c}\NormalTok{(}\ConstantTok{TRUE}\NormalTok{, }\ConstantTok{TRUE}\NormalTok{, }\ConstantTok{FALSE}\NormalTok{))}
\FunctionTok{print}\NormalTok{(al3)}
\end{Highlighting}
\end{Shaded}

\begin{verbatim}
[1] FALSE
\end{verbatim}

\subsubsection{Check missing values -
is.na()}\label{check-missing-values---is.na}

\begin{Shaded}
\begin{Highlighting}[]
\NormalTok{na1 }\OtherTok{\textless{}{-}} \FunctionTok{is.na}\NormalTok{(}\FunctionTok{c}\NormalTok{(}\DecValTok{1}\NormalTok{, }\ConstantTok{NA}\NormalTok{, }\DecValTok{3}\NormalTok{))}
\FunctionTok{print}\NormalTok{(na1)}
\end{Highlighting}
\end{Shaded}

\begin{verbatim}
[1] FALSE  TRUE FALSE
\end{verbatim}

\begin{Shaded}
\begin{Highlighting}[]
\NormalTok{na2 }\OtherTok{\textless{}{-}} \FunctionTok{is.na}\NormalTok{(v1)}
\FunctionTok{print}\NormalTok{(na2)}
\end{Highlighting}
\end{Shaded}

\begin{verbatim}
[1] FALSE FALSE FALSE FALSE FALSE FALSE
\end{verbatim}

\begin{Shaded}
\begin{Highlighting}[]
\NormalTok{na3 }\OtherTok{\textless{}{-}} \FunctionTok{is.na}\NormalTok{(}\FunctionTok{c}\NormalTok{(}\ConstantTok{NA}\NormalTok{, }\ConstantTok{NA}\NormalTok{))}
\FunctionTok{print}\NormalTok{(na3)}
\end{Highlighting}
\end{Shaded}

\begin{verbatim}
[1] TRUE TRUE
\end{verbatim}

\subsubsection{Remove missing values -
na.omit()}\label{remove-missing-values---na.omit}

\begin{Shaded}
\begin{Highlighting}[]
\NormalTok{no1 }\OtherTok{\textless{}{-}} \FunctionTok{na.omit}\NormalTok{(}\FunctionTok{c}\NormalTok{(}\DecValTok{1}\NormalTok{, }\ConstantTok{NA}\NormalTok{, }\DecValTok{3}\NormalTok{))}
\FunctionTok{print}\NormalTok{(no1)}
\end{Highlighting}
\end{Shaded}

\begin{verbatim}
[1] 1 3
attr(,"na.action")
[1] 2
attr(,"class")
[1] "omit"
\end{verbatim}

\begin{Shaded}
\begin{Highlighting}[]
\NormalTok{no2 }\OtherTok{\textless{}{-}} \FunctionTok{na.omit}\NormalTok{(}\FunctionTok{c}\NormalTok{(}\ConstantTok{NA}\NormalTok{, }\DecValTok{2}\NormalTok{, }\ConstantTok{NA}\NormalTok{, }\DecValTok{4}\NormalTok{))}
\FunctionTok{print}\NormalTok{(no2)}
\end{Highlighting}
\end{Shaded}

\begin{verbatim}
[1] 2 4
attr(,"na.action")
[1] 1 3
attr(,"class")
[1] "omit"
\end{verbatim}

\begin{Shaded}
\begin{Highlighting}[]
\NormalTok{no3 }\OtherTok{\textless{}{-}} \FunctionTok{na.omit}\NormalTok{(}\FunctionTok{c}\NormalTok{(}\StringTok{"a"}\NormalTok{, }\ConstantTok{NA}\NormalTok{, }\StringTok{"b"}\NormalTok{))}
\FunctionTok{print}\NormalTok{(no3)}
\end{Highlighting}
\end{Shaded}

\begin{verbatim}
[1] "a" "b"
attr(,"na.action")
[1] 2
attr(,"class")
[1] "omit"
\end{verbatim}

\subsubsection{Replace values -
replace()}\label{replace-values---replace}

\begin{Shaded}
\begin{Highlighting}[]
\NormalTok{rep1 }\OtherTok{\textless{}{-}} \FunctionTok{replace}\NormalTok{(v1, v1 }\SpecialCharTok{\textgreater{}} \DecValTok{4}\NormalTok{, }\DecValTok{0}\NormalTok{)}
\FunctionTok{print}\NormalTok{(rep1)}
\end{Highlighting}
\end{Shaded}

\begin{verbatim}
[1] 1 2 3 4 0 0
\end{verbatim}

\begin{Shaded}
\begin{Highlighting}[]
\NormalTok{rep2 }\OtherTok{\textless{}{-}} \FunctionTok{replace}\NormalTok{(v1, v1 }\SpecialCharTok{==} \DecValTok{3}\NormalTok{, }\DecValTok{99}\NormalTok{)}
\FunctionTok{print}\NormalTok{(rep2)}
\end{Highlighting}
\end{Shaded}

\begin{verbatim}
[1]  1  2 99  4  5  6
\end{verbatim}

\begin{Shaded}
\begin{Highlighting}[]
\NormalTok{rep3 }\OtherTok{\textless{}{-}} \FunctionTok{replace}\NormalTok{(}\FunctionTok{c}\NormalTok{(}\StringTok{"a"}\NormalTok{, }\StringTok{"b"}\NormalTok{, }\StringTok{"c"}\NormalTok{), }\DecValTok{2}\NormalTok{, }\StringTok{"z"}\NormalTok{)}
\FunctionTok{print}\NormalTok{(rep3)}
\end{Highlighting}
\end{Shaded}

\begin{verbatim}
[1] "a" "z" "c"
\end{verbatim}

\subsubsection{Combine strings - paste()}\label{combine-strings---paste}

\begin{Shaded}
\begin{Highlighting}[]
\NormalTok{p1 }\OtherTok{\textless{}{-}} \FunctionTok{paste}\NormalTok{(}\StringTok{"Hello"}\NormalTok{, }\StringTok{"World"}\NormalTok{)}
\FunctionTok{print}\NormalTok{(p1)}
\end{Highlighting}
\end{Shaded}

\begin{verbatim}
[1] "Hello World"
\end{verbatim}

\begin{Shaded}
\begin{Highlighting}[]
\NormalTok{p2 }\OtherTok{\textless{}{-}} \FunctionTok{paste}\NormalTok{(}\StringTok{"Data"}\NormalTok{, }\StringTok{"Science"}\NormalTok{, }\AttributeTok{sep =} \StringTok{"{-}"}\NormalTok{)}
\FunctionTok{print}\NormalTok{(p2)}
\end{Highlighting}
\end{Shaded}

\begin{verbatim}
[1] "Data-Science"
\end{verbatim}

\begin{Shaded}
\begin{Highlighting}[]
\NormalTok{p3 }\OtherTok{\textless{}{-}} \FunctionTok{paste}\NormalTok{(v2, }\AttributeTok{collapse =} \StringTok{", "}\NormalTok{)}
\FunctionTok{print}\NormalTok{(p3)}
\end{Highlighting}
\end{Shaded}

\begin{verbatim}
[1] "Anuj , Tushar, datta "
\end{verbatim}

\subsubsection{Combine strings without space -
paste0()}\label{combine-strings-without-space---paste0}

\begin{Shaded}
\begin{Highlighting}[]
\NormalTok{p01 }\OtherTok{\textless{}{-}} \FunctionTok{paste0}\NormalTok{(}\StringTok{"R"}\NormalTok{, }\StringTok{"Language"}\NormalTok{)}
\FunctionTok{print}\NormalTok{(p01)}
\end{Highlighting}
\end{Shaded}

\begin{verbatim}
[1] "RLanguage"
\end{verbatim}

\begin{Shaded}
\begin{Highlighting}[]
\NormalTok{p02 }\OtherTok{\textless{}{-}} \FunctionTok{paste0}\NormalTok{(}\StringTok{"Roll"}\NormalTok{, }\DecValTok{1}\SpecialCharTok{:}\DecValTok{3}\NormalTok{)}
\FunctionTok{print}\NormalTok{(p02)}
\end{Highlighting}
\end{Shaded}

\begin{verbatim}
[1] "Roll1" "Roll2" "Roll3"
\end{verbatim}

\begin{Shaded}
\begin{Highlighting}[]
\NormalTok{p03 }\OtherTok{\textless{}{-}} \FunctionTok{paste0}\NormalTok{(v2, }\StringTok{"Student"}\NormalTok{)}
\FunctionTok{print}\NormalTok{(p03)}
\end{Highlighting}
\end{Shaded}

\begin{verbatim}
[1] "Anuj Student"  "TusharStudent" "datta Student"
\end{verbatim}

\subsubsection{Convert to uppercase -
toupper()}\label{convert-to-uppercase---toupper}

\begin{Shaded}
\begin{Highlighting}[]
\NormalTok{up1 }\OtherTok{\textless{}{-}} \FunctionTok{toupper}\NormalTok{(}\StringTok{"hello"}\NormalTok{)}
\FunctionTok{print}\NormalTok{(up1)}
\end{Highlighting}
\end{Shaded}

\begin{verbatim}
[1] "HELLO"
\end{verbatim}

\begin{Shaded}
\begin{Highlighting}[]
\NormalTok{up2 }\OtherTok{\textless{}{-}} \FunctionTok{toupper}\NormalTok{(v2)}
\FunctionTok{print}\NormalTok{(up2)}
\end{Highlighting}
\end{Shaded}

\begin{verbatim}
[1] "ANUJ "  "TUSHAR" "DATTA "
\end{verbatim}

\begin{Shaded}
\begin{Highlighting}[]
\NormalTok{up3 }\OtherTok{\textless{}{-}} \FunctionTok{toupper}\NormalTok{(}\FunctionTok{c}\NormalTok{(}\StringTok{"data"}\NormalTok{, }\StringTok{"science"}\NormalTok{))}
\FunctionTok{print}\NormalTok{(up3)}
\end{Highlighting}
\end{Shaded}

\begin{verbatim}
[1] "DATA"    "SCIENCE"
\end{verbatim}

\subsubsection{Convert to lowercase -
tolower()}\label{convert-to-lowercase---tolower}

\begin{Shaded}
\begin{Highlighting}[]
\NormalTok{lo1 }\OtherTok{\textless{}{-}} \FunctionTok{tolower}\NormalTok{(}\StringTok{"HELLO"}\NormalTok{)}
\FunctionTok{print}\NormalTok{(lo1)}
\end{Highlighting}
\end{Shaded}

\begin{verbatim}
[1] "hello"
\end{verbatim}

\begin{Shaded}
\begin{Highlighting}[]
\NormalTok{lo2 }\OtherTok{\textless{}{-}} \FunctionTok{tolower}\NormalTok{(v2)}
\FunctionTok{print}\NormalTok{(lo2)}
\end{Highlighting}
\end{Shaded}

\begin{verbatim}
[1] "anuj "  "tushar" "datta "
\end{verbatim}

\begin{Shaded}
\begin{Highlighting}[]
\NormalTok{lo3 }\OtherTok{\textless{}{-}} \FunctionTok{tolower}\NormalTok{(}\FunctionTok{c}\NormalTok{(}\StringTok{"DATA"}\NormalTok{, }\StringTok{"ML"}\NormalTok{))}
\FunctionTok{print}\NormalTok{(lo3)}
\end{Highlighting}
\end{Shaded}

\begin{verbatim}
[1] "data" "ml"  
\end{verbatim}

\subsubsection{Substring extraction -
substr()}\label{substring-extraction---substr}

\begin{Shaded}
\begin{Highlighting}[]
\NormalTok{sb1 }\OtherTok{\textless{}{-}} \FunctionTok{substr}\NormalTok{(}\StringTok{"datascience"}\NormalTok{, }\DecValTok{1}\NormalTok{, }\DecValTok{4}\NormalTok{)}
\FunctionTok{print}\NormalTok{(sb1)}
\end{Highlighting}
\end{Shaded}

\begin{verbatim}
[1] "data"
\end{verbatim}

\begin{Shaded}
\begin{Highlighting}[]
\NormalTok{sb2 }\OtherTok{\textless{}{-}} \FunctionTok{substr}\NormalTok{(}\StringTok{"machinelearning"}\NormalTok{, }\DecValTok{8}\NormalTok{, }\DecValTok{15}\NormalTok{)}
\FunctionTok{print}\NormalTok{(sb2)}
\end{Highlighting}
\end{Shaded}

\begin{verbatim}
[1] "learning"
\end{verbatim}

\begin{Shaded}
\begin{Highlighting}[]
\NormalTok{sb3 }\OtherTok{\textless{}{-}} \FunctionTok{substr}\NormalTok{(v2, }\DecValTok{1}\NormalTok{, }\DecValTok{3}\NormalTok{)}
\FunctionTok{print}\NormalTok{(sb3)}
\end{Highlighting}
\end{Shaded}

\begin{verbatim}
[1] "Anu" "Tus" "dat"
\end{verbatim}

\subsubsection{Repeat values - rep()}\label{repeat-values---rep}

\begin{Shaded}
\begin{Highlighting}[]
\NormalTok{rp1 }\OtherTok{\textless{}{-}} \FunctionTok{rep}\NormalTok{(}\DecValTok{1}\NormalTok{, }\DecValTok{5}\NormalTok{)}
\FunctionTok{print}\NormalTok{(rp1)}
\end{Highlighting}
\end{Shaded}

\begin{verbatim}
[1] 1 1 1 1 1
\end{verbatim}

\begin{Shaded}
\begin{Highlighting}[]
\NormalTok{rp2 }\OtherTok{\textless{}{-}} \FunctionTok{rep}\NormalTok{(}\FunctionTok{c}\NormalTok{(}\DecValTok{1}\NormalTok{,}\DecValTok{2}\NormalTok{), }\AttributeTok{times =} \DecValTok{3}\NormalTok{)}
\FunctionTok{print}\NormalTok{(rp2)}
\end{Highlighting}
\end{Shaded}

\begin{verbatim}
[1] 1 2 1 2 1 2
\end{verbatim}

\begin{Shaded}
\begin{Highlighting}[]
\NormalTok{rp3 }\OtherTok{\textless{}{-}} \FunctionTok{rep}\NormalTok{(}\FunctionTok{c}\NormalTok{(}\StringTok{"a"}\NormalTok{,}\StringTok{"b"}\NormalTok{), }\AttributeTok{each =} \DecValTok{2}\NormalTok{)}
\FunctionTok{print}\NormalTok{(rp3)}
\end{Highlighting}
\end{Shaded}

\begin{verbatim}
[1] "a" "a" "b" "b"
\end{verbatim}

\subsubsection{Sum of values - sum()}\label{sum-of-values---sum}

\begin{Shaded}
\begin{Highlighting}[]
\NormalTok{sm1 }\OtherTok{\textless{}{-}} \FunctionTok{sum}\NormalTok{(v1)}
\FunctionTok{print}\NormalTok{(sm1)}
\end{Highlighting}
\end{Shaded}

\begin{verbatim}
[1] 21
\end{verbatim}

\begin{Shaded}
\begin{Highlighting}[]
\NormalTok{sm2 }\OtherTok{\textless{}{-}} \FunctionTok{sum}\NormalTok{(}\FunctionTok{c}\NormalTok{(}\DecValTok{1}\NormalTok{, }\DecValTok{2}\NormalTok{, }\DecValTok{3}\NormalTok{, }\ConstantTok{NA}\NormalTok{), }\AttributeTok{na.rm =} \ConstantTok{TRUE}\NormalTok{)}
\FunctionTok{print}\NormalTok{(sm2)}
\end{Highlighting}
\end{Shaded}

\begin{verbatim}
[1] 6
\end{verbatim}

\begin{Shaded}
\begin{Highlighting}[]
\NormalTok{sm3 }\OtherTok{\textless{}{-}} \FunctionTok{sum}\NormalTok{(}\DecValTok{1}\SpecialCharTok{:}\DecValTok{10}\NormalTok{)}
\FunctionTok{print}\NormalTok{(sm3)}
\end{Highlighting}
\end{Shaded}

\begin{verbatim}
[1] 55
\end{verbatim}

\subsubsection{Mean of values - mean()}\label{mean-of-values---mean}

\begin{Shaded}
\begin{Highlighting}[]
\NormalTok{mn1 }\OtherTok{\textless{}{-}} \FunctionTok{mean}\NormalTok{(v1)}
\FunctionTok{print}\NormalTok{(mn1)}
\end{Highlighting}
\end{Shaded}

\begin{verbatim}
[1] 3.5
\end{verbatim}

\begin{Shaded}
\begin{Highlighting}[]
\NormalTok{mn2 }\OtherTok{\textless{}{-}} \FunctionTok{mean}\NormalTok{(}\FunctionTok{c}\NormalTok{(}\DecValTok{1}\NormalTok{, }\DecValTok{2}\NormalTok{, }\DecValTok{3}\NormalTok{, }\ConstantTok{NA}\NormalTok{), }\AttributeTok{na.rm =} \ConstantTok{TRUE}\NormalTok{)}
\FunctionTok{print}\NormalTok{(mn2)}
\end{Highlighting}
\end{Shaded}

\begin{verbatim}
[1] 2
\end{verbatim}

\begin{Shaded}
\begin{Highlighting}[]
\NormalTok{mn3 }\OtherTok{\textless{}{-}} \FunctionTok{mean}\NormalTok{(}\DecValTok{1}\SpecialCharTok{:}\DecValTok{10}\NormalTok{)}
\FunctionTok{print}\NormalTok{(mn3)}
\end{Highlighting}
\end{Shaded}

\begin{verbatim}
[1] 5.5
\end{verbatim}

\subsubsection{Maximum value - max()}\label{maximum-value---max}

\begin{Shaded}
\begin{Highlighting}[]
\NormalTok{mx1 }\OtherTok{\textless{}{-}} \FunctionTok{max}\NormalTok{(v1)}
\FunctionTok{print}\NormalTok{(mx1)}
\end{Highlighting}
\end{Shaded}

\begin{verbatim}
[1] 6
\end{verbatim}

\begin{Shaded}
\begin{Highlighting}[]
\NormalTok{mx2 }\OtherTok{\textless{}{-}} \FunctionTok{max}\NormalTok{(}\FunctionTok{c}\NormalTok{(}\DecValTok{3}\NormalTok{, }\DecValTok{7}\NormalTok{, }\DecValTok{2}\NormalTok{))}
\FunctionTok{print}\NormalTok{(mx2)}
\end{Highlighting}
\end{Shaded}

\begin{verbatim}
[1] 7
\end{verbatim}

\begin{Shaded}
\begin{Highlighting}[]
\NormalTok{mx3 }\OtherTok{\textless{}{-}} \FunctionTok{max}\NormalTok{(}\DecValTok{1}\SpecialCharTok{:}\DecValTok{100}\NormalTok{)}
\FunctionTok{print}\NormalTok{(mx3)}
\end{Highlighting}
\end{Shaded}

\begin{verbatim}
[1] 100
\end{verbatim}

\subsubsection{Minimum value - min()}\label{minimum-value---min}

\begin{Shaded}
\begin{Highlighting}[]
\NormalTok{mnm1 }\OtherTok{\textless{}{-}} \FunctionTok{min}\NormalTok{(v1)}
\FunctionTok{print}\NormalTok{(mnm1)}
\end{Highlighting}
\end{Shaded}

\begin{verbatim}
[1] 1
\end{verbatim}

\begin{Shaded}
\begin{Highlighting}[]
\NormalTok{mnm2 }\OtherTok{\textless{}{-}} \FunctionTok{min}\NormalTok{(}\FunctionTok{c}\NormalTok{(}\DecValTok{3}\NormalTok{, }\DecValTok{7}\NormalTok{, }\DecValTok{2}\NormalTok{))}
\FunctionTok{print}\NormalTok{(mnm2)}
\end{Highlighting}
\end{Shaded}

\begin{verbatim}
[1] 2
\end{verbatim}

\begin{Shaded}
\begin{Highlighting}[]
\NormalTok{mnm3 }\OtherTok{\textless{}{-}} \FunctionTok{min}\NormalTok{(}\DecValTok{1}\SpecialCharTok{:}\DecValTok{100}\NormalTok{)}
\FunctionTok{print}\NormalTok{(mnm3) }
\end{Highlighting}
\end{Shaded}

\begin{verbatim}
[1] 1
\end{verbatim}

\subsection{DATA MANIPULATION IN R}\label{data-manipulation-in-r}

\subsubsection{Create a data frame -
data.frame()}\label{create-a-data-frame---data.frame}

\begin{Shaded}
\begin{Highlighting}[]
\CommentTok{\# Simple data frame}
\NormalTok{df1 }\OtherTok{\textless{}{-}} \FunctionTok{data.frame}\NormalTok{(}
  \AttributeTok{id =} \FunctionTok{c}\NormalTok{(}\DecValTok{1}\NormalTok{,}\DecValTok{2}\NormalTok{,}\DecValTok{3}\NormalTok{),}
  \AttributeTok{name =} \FunctionTok{c}\NormalTok{(}\StringTok{"A"}\NormalTok{,}\StringTok{"B"}\NormalTok{,}\StringTok{"C"}\NormalTok{),}
  \AttributeTok{marks =} \FunctionTok{c}\NormalTok{(}\DecValTok{70}\NormalTok{,}\DecValTok{85}\NormalTok{,}\DecValTok{90}\NormalTok{)}
\NormalTok{)}
\FunctionTok{print}\NormalTok{(df1)}
\end{Highlighting}
\end{Shaded}

\begin{verbatim}
  id name marks
1  1    A    70
2  2    B    85
3  3    C    90
\end{verbatim}

\begin{Shaded}
\begin{Highlighting}[]
\CommentTok{\# Data frame without factors}
\NormalTok{df2 }\OtherTok{\textless{}{-}} \FunctionTok{data.frame}\NormalTok{(}
  \AttributeTok{city =} \FunctionTok{c}\NormalTok{(}\StringTok{"Delhi"}\NormalTok{,}\StringTok{"Mumbai"}\NormalTok{,}\StringTok{"Pune"}\NormalTok{),}
  \AttributeTok{temp =} \FunctionTok{c}\NormalTok{(}\DecValTok{35}\NormalTok{,}\DecValTok{30}\NormalTok{,}\DecValTok{28}\NormalTok{),}
  \AttributeTok{stringsAsFactors =} \ConstantTok{FALSE}
\NormalTok{)}
\FunctionTok{print}\NormalTok{(df2)}
\end{Highlighting}
\end{Shaded}

\begin{verbatim}
    city temp
1  Delhi   35
2 Mumbai   30
3   Pune   28
\end{verbatim}

\begin{Shaded}
\begin{Highlighting}[]
\CommentTok{\# Mixed data types}
\NormalTok{df3 }\OtherTok{\textless{}{-}} \FunctionTok{data.frame}\NormalTok{(}
  \AttributeTok{roll =} \DecValTok{1}\SpecialCharTok{:}\DecValTok{3}\NormalTok{,}
  \AttributeTok{pass =} \FunctionTok{c}\NormalTok{(}\ConstantTok{TRUE}\NormalTok{, }\ConstantTok{FALSE}\NormalTok{, }\ConstantTok{TRUE}\NormalTok{)}
\NormalTok{)}
\FunctionTok{print}\NormalTok{(df3)}
\end{Highlighting}
\end{Shaded}

\begin{verbatim}
  roll  pass
1    1  TRUE
2    2 FALSE
3    3  TRUE
\end{verbatim}

\subsubsection{Structure and summary - str(),
summary()}\label{structure-and-summary---str-summary}

\begin{Shaded}
\begin{Highlighting}[]
\FunctionTok{str}\NormalTok{(df1)}
\end{Highlighting}
\end{Shaded}

\begin{verbatim}
'data.frame':   3 obs. of  3 variables:
 $ id   : num  1 2 3
 $ name : chr  "A" "B" "C"
 $ marks: num  70 85 90
\end{verbatim}

\begin{Shaded}
\begin{Highlighting}[]
\FunctionTok{summary}\NormalTok{(df1)}
\end{Highlighting}
\end{Shaded}

\begin{verbatim}
       id          name               marks      
 Min.   :1.0   Length:3           Min.   :70.00  
 1st Qu.:1.5   Class :character   1st Qu.:77.50  
 Median :2.0   Mode  :character   Median :85.00  
 Mean   :2.0                      Mean   :81.67  
 3rd Qu.:2.5                      3rd Qu.:87.50  
 Max.   :3.0                      Max.   :90.00  
\end{verbatim}

\begin{Shaded}
\begin{Highlighting}[]
\FunctionTok{summary}\NormalTok{(df2)}
\end{Highlighting}
\end{Shaded}

\begin{verbatim}
     city                temp     
 Length:3           Min.   :28.0  
 Class :character   1st Qu.:29.0  
 Mode  :character   Median :30.0  
                    Mean   :31.0  
                    3rd Qu.:32.5  
                    Max.   :35.0  
\end{verbatim}

\subsubsection{Access columns - \$, {[} {]}}\label{access-columns--}

\begin{Shaded}
\begin{Highlighting}[]
\NormalTok{df1}\SpecialCharTok{$}\NormalTok{marks}
\end{Highlighting}
\end{Shaded}

\begin{verbatim}
[1] 70 85 90
\end{verbatim}

\begin{Shaded}
\begin{Highlighting}[]
\NormalTok{df1[,}\DecValTok{2}\NormalTok{]}
\end{Highlighting}
\end{Shaded}

\begin{verbatim}
[1] "A" "B" "C"
\end{verbatim}

\begin{Shaded}
\begin{Highlighting}[]
\NormalTok{df1[, }\FunctionTok{c}\NormalTok{(}\StringTok{"name"}\NormalTok{,}\StringTok{"marks"}\NormalTok{)]}
\end{Highlighting}
\end{Shaded}

\begin{verbatim}
  name marks
1    A    70
2    B    85
3    C    90
\end{verbatim}

\subsubsection{Add new columns}\label{add-new-columns}

\begin{Shaded}
\begin{Highlighting}[]
\NormalTok{df1}\SpecialCharTok{$}\NormalTok{result }\OtherTok{\textless{}{-}} \FunctionTok{c}\NormalTok{(}\StringTok{"Pass"}\NormalTok{,}\StringTok{"Pass"}\NormalTok{,}\StringTok{"Pass"}\NormalTok{)}
\FunctionTok{print}\NormalTok{(df1)}
\end{Highlighting}
\end{Shaded}

\begin{verbatim}
  id name marks result
1  1    A    70   Pass
2  2    B    85   Pass
3  3    C    90   Pass
\end{verbatim}

\begin{Shaded}
\begin{Highlighting}[]
\NormalTok{df1}\SpecialCharTok{$}\NormalTok{grade }\OtherTok{\textless{}{-}} \FunctionTok{ifelse}\NormalTok{(df1}\SpecialCharTok{$}\NormalTok{marks }\SpecialCharTok{\textgreater{}=} \DecValTok{80}\NormalTok{, }\StringTok{"A"}\NormalTok{, }\StringTok{"B"}\NormalTok{)}
\FunctionTok{print}\NormalTok{(df1)}
\end{Highlighting}
\end{Shaded}

\begin{verbatim}
  id name marks result grade
1  1    A    70   Pass     B
2  2    B    85   Pass     A
3  3    C    90   Pass     A
\end{verbatim}

\begin{Shaded}
\begin{Highlighting}[]
\NormalTok{df1}\SpecialCharTok{$}\NormalTok{distinction }\OtherTok{\textless{}{-}}\NormalTok{ df1}\SpecialCharTok{$}\NormalTok{marks }\SpecialCharTok{\textgreater{}} \DecValTok{85}
\FunctionTok{print}\NormalTok{(df1)}
\end{Highlighting}
\end{Shaded}

\begin{verbatim}
  id name marks result grade distinction
1  1    A    70   Pass     B       FALSE
2  2    B    85   Pass     A       FALSE
3  3    C    90   Pass     A        TRUE
\end{verbatim}

\subsubsection{Remove columns}\label{remove-columns}

\begin{Shaded}
\begin{Highlighting}[]
\NormalTok{df1}\SpecialCharTok{$}\NormalTok{result }\OtherTok{\textless{}{-}} \ConstantTok{NULL}
\FunctionTok{print}\NormalTok{(df1)}
\end{Highlighting}
\end{Shaded}

\begin{verbatim}
  id name marks grade distinction
1  1    A    70     B       FALSE
2  2    B    85     A       FALSE
3  3    C    90     A        TRUE
\end{verbatim}

\begin{Shaded}
\begin{Highlighting}[]
\NormalTok{df1 }\OtherTok{\textless{}{-}}\NormalTok{ df1[, }\SpecialCharTok{{-}}\DecValTok{4}\NormalTok{]}
\FunctionTok{print}\NormalTok{(df1)}
\end{Highlighting}
\end{Shaded}

\begin{verbatim}
  id name marks distinction
1  1    A    70       FALSE
2  2    B    85       FALSE
3  3    C    90        TRUE
\end{verbatim}

\begin{Shaded}
\begin{Highlighting}[]
\NormalTok{df1 }\OtherTok{\textless{}{-}}\NormalTok{ df1[, }\SpecialCharTok{!}\NormalTok{(}\FunctionTok{names}\NormalTok{(df1) }\SpecialCharTok{\%in\%} \FunctionTok{c}\NormalTok{(}\StringTok{"name"}\NormalTok{))]}
\FunctionTok{print}\NormalTok{(df1)}
\end{Highlighting}
\end{Shaded}

\begin{verbatim}
  id marks distinction
1  1    70       FALSE
2  2    85       FALSE
3  3    90        TRUE
\end{verbatim}

\subsubsection{Filter rows - subset()}\label{filter-rows---subset}

\begin{Shaded}
\begin{Highlighting}[]
\FunctionTok{subset}\NormalTok{(df1, marks }\SpecialCharTok{\textgreater{}=} \DecValTok{80}\NormalTok{)}
\end{Highlighting}
\end{Shaded}

\begin{verbatim}
  id marks distinction
2  2    85       FALSE
3  3    90        TRUE
\end{verbatim}

\subsubsection{Conditional logic -
ifelse()}\label{conditional-logic---ifelse}

\begin{Shaded}
\begin{Highlighting}[]
\NormalTok{df1}\SpecialCharTok{$}\NormalTok{status }\OtherTok{\textless{}{-}} \FunctionTok{ifelse}\NormalTok{(df1}\SpecialCharTok{$}\NormalTok{marks }\SpecialCharTok{\textgreater{}=} \DecValTok{75}\NormalTok{, }\StringTok{"Pass"}\NormalTok{, }\StringTok{"Fail"}\NormalTok{)}
\FunctionTok{print}\NormalTok{(df1)}
\end{Highlighting}
\end{Shaded}

\begin{verbatim}
  id marks distinction status
1  1    70       FALSE   Fail
2  2    85       FALSE   Pass
3  3    90        TRUE   Pass
\end{verbatim}

\begin{Shaded}
\begin{Highlighting}[]
\NormalTok{df1}\SpecialCharTok{$}\NormalTok{grade2 }\OtherTok{\textless{}{-}} \FunctionTok{ifelse}\NormalTok{(df1}\SpecialCharTok{$}\NormalTok{marks }\SpecialCharTok{\textgreater{}=} \DecValTok{85}\NormalTok{, }\StringTok{"A"}\NormalTok{,}
                     \FunctionTok{ifelse}\NormalTok{(df1}\SpecialCharTok{$}\NormalTok{marks }\SpecialCharTok{\textgreater{}=} \DecValTok{70}\NormalTok{, }\StringTok{"B"}\NormalTok{, }\StringTok{"C"}\NormalTok{))}
\FunctionTok{print}\NormalTok{(df1)}
\end{Highlighting}
\end{Shaded}

\begin{verbatim}
  id marks distinction status grade2
1  1    70       FALSE   Fail      B
2  2    85       FALSE   Pass      A
3  3    90        TRUE   Pass      A
\end{verbatim}

\begin{Shaded}
\begin{Highlighting}[]
\FunctionTok{ifelse}\NormalTok{(df1}\SpecialCharTok{$}\NormalTok{marks }\SpecialCharTok{\textgreater{}} \DecValTok{80}\NormalTok{, }\ConstantTok{TRUE}\NormalTok{, }\ConstantTok{FALSE}\NormalTok{)}
\end{Highlighting}
\end{Shaded}

\begin{verbatim}
[1] FALSE  TRUE  TRUE
\end{verbatim}

\subsubsection{Sorting data - order()}\label{sorting-data---order}

\begin{Shaded}
\begin{Highlighting}[]
\NormalTok{df1[}\FunctionTok{order}\NormalTok{(df1}\SpecialCharTok{$}\NormalTok{marks), ]}
\end{Highlighting}
\end{Shaded}

\begin{verbatim}
  id marks distinction status grade2
1  1    70       FALSE   Fail      B
2  2    85       FALSE   Pass      A
3  3    90        TRUE   Pass      A
\end{verbatim}

\begin{Shaded}
\begin{Highlighting}[]
\NormalTok{df1[}\FunctionTok{order}\NormalTok{(df1}\SpecialCharTok{$}\NormalTok{marks, }\AttributeTok{decreasing =} \ConstantTok{TRUE}\NormalTok{), ]}
\end{Highlighting}
\end{Shaded}

\begin{verbatim}
  id marks distinction status grade2
3  3    90        TRUE   Pass      A
2  2    85       FALSE   Pass      A
1  1    70       FALSE   Fail      B
\end{verbatim}

\begin{Shaded}
\begin{Highlighting}[]
\NormalTok{df1[}\FunctionTok{order}\NormalTok{(df1}\SpecialCharTok{$}\NormalTok{grade, df1}\SpecialCharTok{$}\NormalTok{marks), ]}
\end{Highlighting}
\end{Shaded}

\begin{verbatim}
  id marks distinction status grade2
2  2    85       FALSE   Pass      A
3  3    90        TRUE   Pass      A
1  1    70       FALSE   Fail      B
\end{verbatim}

\subsubsection{Aggregation - aggregate()}\label{aggregation---aggregate}

\begin{Shaded}
\begin{Highlighting}[]
\FunctionTok{aggregate}\NormalTok{(marks }\SpecialCharTok{\textasciitilde{}}\NormalTok{ status, }\AttributeTok{data =}\NormalTok{ df1, mean)}
\end{Highlighting}
\end{Shaded}

\begin{verbatim}
  status marks
1   Fail  70.0
2   Pass  87.5
\end{verbatim}

\begin{Shaded}
\begin{Highlighting}[]
\FunctionTok{aggregate}\NormalTok{(id }\SpecialCharTok{\textasciitilde{}}\NormalTok{ status, }\AttributeTok{data =}\NormalTok{ df1, length)}
\end{Highlighting}
\end{Shaded}

\begin{verbatim}
  status id
1   Fail  1
2   Pass  2
\end{verbatim}

\begin{Shaded}
\begin{Highlighting}[]
\FunctionTok{aggregate}\NormalTok{(marks }\SpecialCharTok{\textasciitilde{}}\NormalTok{ status, }\AttributeTok{data =}\NormalTok{ df1, sum)}
\end{Highlighting}
\end{Shaded}

\begin{verbatim}
  status marks
1   Fail    70
2   Pass   175
\end{verbatim}

\subsubsection{Merge data frames -
merge()}\label{merge-data-frames---merge}

\begin{Shaded}
\begin{Highlighting}[]
\NormalTok{df4 }\OtherTok{\textless{}{-}} \FunctionTok{data.frame}\NormalTok{(}
  \AttributeTok{id =} \FunctionTok{c}\NormalTok{(}\DecValTok{1}\NormalTok{,}\DecValTok{2}\NormalTok{,}\DecValTok{3}\NormalTok{),}
  \AttributeTok{age =} \FunctionTok{c}\NormalTok{(}\DecValTok{20}\NormalTok{,}\DecValTok{21}\NormalTok{,}\DecValTok{22}\NormalTok{)}
\NormalTok{)}

\FunctionTok{merge}\NormalTok{(df1, df4, }\AttributeTok{by =} \StringTok{"id"}\NormalTok{)}
\end{Highlighting}
\end{Shaded}

\begin{verbatim}
  id marks distinction status grade2 age
1  1    70       FALSE   Fail      B  20
2  2    85       FALSE   Pass      A  21
3  3    90        TRUE   Pass      A  22
\end{verbatim}

\begin{Shaded}
\begin{Highlighting}[]
\FunctionTok{merge}\NormalTok{(df1, df4, }\AttributeTok{by =} \StringTok{"id"}\NormalTok{, }\AttributeTok{all.x =} \ConstantTok{TRUE}\NormalTok{)}
\end{Highlighting}
\end{Shaded}

\begin{verbatim}
  id marks distinction status grade2 age
1  1    70       FALSE   Fail      B  20
2  2    85       FALSE   Pass      A  21
3  3    90        TRUE   Pass      A  22
\end{verbatim}

\begin{Shaded}
\begin{Highlighting}[]
\FunctionTok{merge}\NormalTok{(df1, df4, }\AttributeTok{by =} \StringTok{"id"}\NormalTok{, }\AttributeTok{all.y =} \ConstantTok{TRUE}\NormalTok{)}
\end{Highlighting}
\end{Shaded}

\begin{verbatim}
  id marks distinction status grade2 age
1  1    70       FALSE   Fail      B  20
2  2    85       FALSE   Pass      A  21
3  3    90        TRUE   Pass      A  22
\end{verbatim}

\subsubsection{Apply family - apply(), lapply(),
sapply()}\label{apply-family---apply-lapply-sapply}

\begin{Shaded}
\begin{Highlighting}[]
\FunctionTok{apply}\NormalTok{(df4[, }\StringTok{"age"}\NormalTok{, }\AttributeTok{drop =} \ConstantTok{FALSE}\NormalTok{], }\DecValTok{1}\NormalTok{, mean)}
\end{Highlighting}
\end{Shaded}

\begin{verbatim}
[1] 20 21 22
\end{verbatim}

\begin{Shaded}
\begin{Highlighting}[]
\FunctionTok{apply}\NormalTok{(df4[, }\StringTok{"age"}\NormalTok{, }\AttributeTok{drop =} \ConstantTok{FALSE}\NormalTok{], }\DecValTok{2}\NormalTok{, mean)}
\end{Highlighting}
\end{Shaded}

\begin{verbatim}
age 
 21 
\end{verbatim}

\begin{Shaded}
\begin{Highlighting}[]
\FunctionTok{lapply}\NormalTok{(df4, class)}
\end{Highlighting}
\end{Shaded}

\begin{verbatim}
$id
[1] "numeric"

$age
[1] "numeric"
\end{verbatim}

\begin{Shaded}
\begin{Highlighting}[]
\FunctionTok{sapply}\NormalTok{(df4, class)}
\end{Highlighting}
\end{Shaded}

\begin{verbatim}
       id       age 
"numeric" "numeric" 
\end{verbatim}

\subsubsection{Handling missing values}\label{handling-missing-values}

\begin{Shaded}
\begin{Highlighting}[]
\NormalTok{vec1 }\OtherTok{\textless{}{-}} \FunctionTok{c}\NormalTok{(}\DecValTok{1}\NormalTok{, }\ConstantTok{NA}\NormalTok{, }\DecValTok{3}\NormalTok{, }\ConstantTok{NA}\NormalTok{)}
\FunctionTok{is.na}\NormalTok{(vec1)}
\end{Highlighting}
\end{Shaded}

\begin{verbatim}
[1] FALSE  TRUE FALSE  TRUE
\end{verbatim}

\begin{Shaded}
\begin{Highlighting}[]
\FunctionTok{na.omit}\NormalTok{(vec1)}
\end{Highlighting}
\end{Shaded}

\begin{verbatim}
[1] 1 3
attr(,"na.action")
[1] 2 4
attr(,"class")
[1] "omit"
\end{verbatim}

\begin{Shaded}
\begin{Highlighting}[]
\FunctionTok{sum}\NormalTok{(vec1, }\AttributeTok{na.rm =} \ConstantTok{TRUE}\NormalTok{)}
\end{Highlighting}
\end{Shaded}

\begin{verbatim}
[1] 4
\end{verbatim}

\subsubsection{Factors and levels}\label{factors-and-levels}

\begin{Shaded}
\begin{Highlighting}[]
\NormalTok{df1}\SpecialCharTok{$}\NormalTok{grade }\OtherTok{\textless{}{-}} \FunctionTok{factor}\NormalTok{(df1}\SpecialCharTok{$}\NormalTok{grade)}
\FunctionTok{levels}\NormalTok{(df1}\SpecialCharTok{$}\NormalTok{grade)}
\end{Highlighting}
\end{Shaded}

\begin{verbatim}
[1] "A" "B"
\end{verbatim}

\begin{Shaded}
\begin{Highlighting}[]
\FunctionTok{nlevels}\NormalTok{(df1}\SpecialCharTok{$}\NormalTok{grade)}
\end{Highlighting}
\end{Shaded}

\begin{verbatim}
[1] 2
\end{verbatim}

\begin{Shaded}
\begin{Highlighting}[]
\FunctionTok{as.character}\NormalTok{(df1}\SpecialCharTok{$}\NormalTok{grade)}
\end{Highlighting}
\end{Shaded}

\begin{verbatim}
[1] "B" "A" "A"
\end{verbatim}

\subsubsection{Rename columns}\label{rename-columns}

\begin{Shaded}
\begin{Highlighting}[]
\FunctionTok{colnames}\NormalTok{(df1)[}\FunctionTok{colnames}\NormalTok{(df1) }\SpecialCharTok{==} \StringTok{"marks"}\NormalTok{] }\OtherTok{\textless{}{-}} \StringTok{"score"}
\FunctionTok{print}\NormalTok{(df1)}
\end{Highlighting}
\end{Shaded}

\begin{verbatim}
  id score distinction status grade2 grade
1  1    70       FALSE   Fail      B     B
2  2    85       FALSE   Pass      A     A
3  3    90        TRUE   Pass      A     A
\end{verbatim}

\subsubsection{Sample rows}\label{sample-rows}

\begin{Shaded}
\begin{Highlighting}[]
\FunctionTok{sample}\NormalTok{(df1}\SpecialCharTok{$}\NormalTok{id)}
\end{Highlighting}
\end{Shaded}

\begin{verbatim}
[1] 2 1 3
\end{verbatim}

\begin{Shaded}
\begin{Highlighting}[]
\FunctionTok{sample}\NormalTok{(df1}\SpecialCharTok{$}\NormalTok{id, }\DecValTok{2}\NormalTok{)}
\end{Highlighting}
\end{Shaded}

\begin{verbatim}
[1] 2 1
\end{verbatim}

\begin{Shaded}
\begin{Highlighting}[]
\FunctionTok{sample}\NormalTok{(df1}\SpecialCharTok{$}\NormalTok{id, }\AttributeTok{replace =} \ConstantTok{TRUE}\NormalTok{)}
\end{Highlighting}
\end{Shaded}

\begin{verbatim}
[1] 2 2 1
\end{verbatim}

\subsubsection{Bind data frames - rbind(),
cbind()}\label{bind-data-frames---rbind-cbind}

\begin{Shaded}
\begin{Highlighting}[]
\NormalTok{df5 }\OtherTok{\textless{}{-}} \FunctionTok{data.frame}\NormalTok{(}
  \AttributeTok{id =} \FunctionTok{c}\NormalTok{(}\DecValTok{4}\NormalTok{,}\DecValTok{5}\NormalTok{),}
  \AttributeTok{score =} \FunctionTok{c}\NormalTok{(}\DecValTok{60}\NormalTok{,}\DecValTok{95}\NormalTok{),}
  \AttributeTok{grade =} \FunctionTok{c}\NormalTok{(}\StringTok{"C"}\NormalTok{,}\StringTok{"A"}\NormalTok{),}
  \AttributeTok{distinction =} \FunctionTok{c}\NormalTok{(}\ConstantTok{FALSE}\NormalTok{, }\ConstantTok{TRUE}\NormalTok{),}
  \AttributeTok{status =} \FunctionTok{c}\NormalTok{(}\StringTok{"Fail"}\NormalTok{,}\StringTok{"Pass"}\NormalTok{),}
  \AttributeTok{grade2 =} \FunctionTok{c}\NormalTok{(}\StringTok{"C"}\NormalTok{,}\StringTok{"A"}\NormalTok{)}
\NormalTok{)}

\FunctionTok{rbind}\NormalTok{(df1, df5)}
\end{Highlighting}
\end{Shaded}

\begin{verbatim}
  id score distinction status grade2 grade
1  1    70       FALSE   Fail      B     B
2  2    85       FALSE   Pass      A     A
3  3    90        TRUE   Pass      A     A
4  4    60       FALSE   Fail      C     C
5  5    95        TRUE   Pass      A     A
\end{verbatim}

\begin{Shaded}
\begin{Highlighting}[]
\FunctionTok{cbind}\NormalTok{(df4, }\AttributeTok{gender =} \FunctionTok{c}\NormalTok{(}\StringTok{"M"}\NormalTok{,}\StringTok{"F"}\NormalTok{,}\StringTok{"M"}\NormalTok{))}
\end{Highlighting}
\end{Shaded}

\begin{verbatim}
  id age gender
1  1  20      M
2  2  21      F
3  3  22      M
\end{verbatim}

\subsubsection{Basic transformations}\label{basic-transformations}

\begin{Shaded}
\begin{Highlighting}[]
\NormalTok{df1}\SpecialCharTok{$}\NormalTok{score\_scaled }\OtherTok{\textless{}{-}}\NormalTok{ df1}\SpecialCharTok{$}\NormalTok{score }\SpecialCharTok{/} \FunctionTok{max}\NormalTok{(df1}\SpecialCharTok{$}\NormalTok{score)}
\FunctionTok{print}\NormalTok{(df1)}
\end{Highlighting}
\end{Shaded}

\begin{verbatim}
  id score distinction status grade2 grade score_scaled
1  1    70       FALSE   Fail      B     B    0.7777778
2  2    85       FALSE   Pass      A     A    0.9444444
3  3    90        TRUE   Pass      A     A    1.0000000
\end{verbatim}

\begin{Shaded}
\begin{Highlighting}[]
\FunctionTok{log}\NormalTok{(df1}\SpecialCharTok{$}\NormalTok{score)}
\end{Highlighting}
\end{Shaded}

\begin{verbatim}
[1] 4.248495 4.442651 4.499810
\end{verbatim}

\begin{Shaded}
\begin{Highlighting}[]
\FunctionTok{sqrt}\NormalTok{(df1}\SpecialCharTok{$}\NormalTok{score)}
\end{Highlighting}
\end{Shaded}

\begin{verbatim}
[1] 8.366600 9.219544 9.486833
\end{verbatim}

\subsubsection{Head and tail}\label{head-and-tail}

\begin{Shaded}
\begin{Highlighting}[]
\FunctionTok{head}\NormalTok{(df1)}
\end{Highlighting}
\end{Shaded}

\begin{verbatim}
  id score distinction status grade2 grade score_scaled
1  1    70       FALSE   Fail      B     B    0.7777778
2  2    85       FALSE   Pass      A     A    0.9444444
3  3    90        TRUE   Pass      A     A    1.0000000
\end{verbatim}

\begin{Shaded}
\begin{Highlighting}[]
\FunctionTok{tail}\NormalTok{(df1)}
\end{Highlighting}
\end{Shaded}

\begin{verbatim}
  id score distinction status grade2 grade score_scaled
1  1    70       FALSE   Fail      B     B    0.7777778
2  2    85       FALSE   Pass      A     A    0.9444444
3  3    90        TRUE   Pass      A     A    1.0000000
\end{verbatim}

\subsection{Control Flow \& Programming}\label{control-flow-programming}

\paragraph{ifelse() -- Vectorized if-else (works on each
element)}\label{ifelse-vectorized-if-else-works-on-each-element}

\begin{Shaded}
\begin{Highlighting}[]
\CommentTok{\#   Check if numbers are even or odd}
\NormalTok{result1 }\OtherTok{\textless{}{-}} \FunctionTok{ifelse}\NormalTok{(}\FunctionTok{c}\NormalTok{(}\DecValTok{1}\NormalTok{, }\DecValTok{2}\NormalTok{, }\DecValTok{3}\NormalTok{, }\DecValTok{4}\NormalTok{) }\SpecialCharTok{\%\%} \DecValTok{2} \SpecialCharTok{==} \DecValTok{0}\NormalTok{, }\StringTok{"even"}\NormalTok{, }\StringTok{"odd"}\NormalTok{)}
\FunctionTok{print}\NormalTok{(result1)}
\end{Highlighting}
\end{Shaded}

\begin{verbatim}
[1] "odd"  "even" "odd"  "even"
\end{verbatim}

\begin{Shaded}
\begin{Highlighting}[]
\CommentTok{\#   Mark values above 5 as "high", others "low"}
\NormalTok{result2 }\OtherTok{\textless{}{-}} \FunctionTok{ifelse}\NormalTok{(}\FunctionTok{c}\NormalTok{(}\DecValTok{3}\NormalTok{, }\DecValTok{6}\NormalTok{, }\DecValTok{9}\NormalTok{), }\StringTok{"high"}\NormalTok{, }\StringTok{"low"}\NormalTok{)}
\FunctionTok{print}\NormalTok{(result2)}
\end{Highlighting}
\end{Shaded}

\begin{verbatim}
[1] "high" "high" "high"
\end{verbatim}

\begin{Shaded}
\begin{Highlighting}[]
\CommentTok{\#   Simple if{-}else statement}
\NormalTok{x }\OtherTok{\textless{}{-}} \DecValTok{10}
\ControlFlowTok{if}\NormalTok{ (x }\SpecialCharTok{\textgreater{}} \DecValTok{5}\NormalTok{) \{}
  \FunctionTok{print}\NormalTok{(}\StringTok{"x is greater than 5"}\NormalTok{)}
\NormalTok{\} }\ControlFlowTok{else}\NormalTok{ \{}
  \FunctionTok{print}\NormalTok{(}\StringTok{"x is not greater than 5"}\NormalTok{)}
\NormalTok{\}}
\end{Highlighting}
\end{Shaded}

\begin{verbatim}
[1] "x is greater than 5"
\end{verbatim}

\begin{Shaded}
\begin{Highlighting}[]
\CommentTok{\#   Using if{-}else inside a function}
\NormalTok{check\_number }\OtherTok{\textless{}{-}} \ControlFlowTok{function}\NormalTok{(num) \{}
  \ControlFlowTok{if}\NormalTok{ (num }\SpecialCharTok{\textless{}} \DecValTok{0}\NormalTok{) \{}
    \FunctionTok{return}\NormalTok{(}\StringTok{"Negative"}\NormalTok{)}
\NormalTok{  \} }\ControlFlowTok{else}\NormalTok{ \{}
    \FunctionTok{return}\NormalTok{(}\StringTok{"Non{-}negative"}\NormalTok{)}
\NormalTok{  \}}
\NormalTok{\}}
\FunctionTok{print}\NormalTok{(}\FunctionTok{check\_number}\NormalTok{(}\SpecialCharTok{{-}}\DecValTok{3}\NormalTok{))}
\end{Highlighting}
\end{Shaded}

\begin{verbatim}
[1] "Negative"
\end{verbatim}

\begin{Shaded}
\begin{Highlighting}[]
\FunctionTok{print}\NormalTok{(}\FunctionTok{check\_number}\NormalTok{(}\DecValTok{4}\NormalTok{))}
\end{Highlighting}
\end{Shaded}

\begin{verbatim}
[1] "Non-negative"
\end{verbatim}

\begin{Shaded}
\begin{Highlighting}[]
\CommentTok{\#   Print numbers 1 to 5}
\ControlFlowTok{for}\NormalTok{ (i }\ControlFlowTok{in} \DecValTok{1}\SpecialCharTok{:}\DecValTok{5}\NormalTok{) \{}
  \FunctionTok{print}\NormalTok{(i)}
\NormalTok{\}}
\end{Highlighting}
\end{Shaded}

\begin{verbatim}
[1] 1
[1] 2
[1] 3
[1] 4
[1] 5
\end{verbatim}

\begin{Shaded}
\begin{Highlighting}[]
\CommentTok{\#   Loop over elements in a vector}
\NormalTok{vec\_loop }\OtherTok{\textless{}{-}} \FunctionTok{c}\NormalTok{(}\StringTok{"a"}\NormalTok{, }\StringTok{"b"}\NormalTok{, }\StringTok{"c"}\NormalTok{)}
\ControlFlowTok{for}\NormalTok{ (letter }\ControlFlowTok{in}\NormalTok{ vec\_loop) \{}
  \FunctionTok{print}\NormalTok{(letter)}
\NormalTok{\}}
\end{Highlighting}
\end{Shaded}

\begin{verbatim}
[1] "a"
[1] "b"
[1] "c"
\end{verbatim}

\begin{Shaded}
\begin{Highlighting}[]
\CommentTok{\#   Count from 1 until 5}
\NormalTok{counter }\OtherTok{\textless{}{-}} \DecValTok{1}
\ControlFlowTok{while}\NormalTok{ (counter }\SpecialCharTok{\textless{}=} \DecValTok{5}\NormalTok{) \{}
  \FunctionTok{print}\NormalTok{(counter)}
\NormalTok{  counter }\OtherTok{\textless{}{-}}\NormalTok{ counter }\SpecialCharTok{+} \DecValTok{1}
\NormalTok{\}}
\end{Highlighting}
\end{Shaded}

\begin{verbatim}
[1] 1
[1] 2
[1] 3
[1] 4
[1] 5
\end{verbatim}

\begin{Shaded}
\begin{Highlighting}[]
\CommentTok{\#   Decrease a value until it is less than 0}
\NormalTok{val }\OtherTok{\textless{}{-}} \DecValTok{3}
\ControlFlowTok{while}\NormalTok{ (val }\SpecialCharTok{\textgreater{}=} \DecValTok{0}\NormalTok{) \{}
  \FunctionTok{print}\NormalTok{(val)}
\NormalTok{  val }\OtherTok{\textless{}{-}}\NormalTok{ val }\SpecialCharTok{{-}} \DecValTok{1}
\NormalTok{\}}
\end{Highlighting}
\end{Shaded}

\begin{verbatim}
[1] 3
[1] 2
[1] 1
[1] 0
\end{verbatim}

\begin{Shaded}
\begin{Highlighting}[]
\CommentTok{\#   Use repeat loop to count until 3}
\NormalTok{count }\OtherTok{\textless{}{-}} \DecValTok{1}
\ControlFlowTok{repeat}\NormalTok{ \{}
  \FunctionTok{print}\NormalTok{(count)}
\NormalTok{  count }\OtherTok{\textless{}{-}}\NormalTok{ count }\SpecialCharTok{+} \DecValTok{1}
  \ControlFlowTok{if}\NormalTok{ (count }\SpecialCharTok{\textgreater{}} \DecValTok{3}\NormalTok{) }\ControlFlowTok{break}
\NormalTok{\}}
\end{Highlighting}
\end{Shaded}

\begin{verbatim}
[1] 1
[1] 2
[1] 3
\end{verbatim}

\begin{Shaded}
\begin{Highlighting}[]
\CommentTok{\#   Another repeat example with a condition}
\NormalTok{i }\OtherTok{\textless{}{-}} \DecValTok{10}
\ControlFlowTok{repeat}\NormalTok{ \{}
\NormalTok{  i }\OtherTok{\textless{}{-}}\NormalTok{ i }\SpecialCharTok{{-}} \DecValTok{2}
  \FunctionTok{print}\NormalTok{(i)}
  \ControlFlowTok{if}\NormalTok{ (i }\SpecialCharTok{\textless{}=} \DecValTok{4}\NormalTok{) }\ControlFlowTok{break}
\NormalTok{\}}
\end{Highlighting}
\end{Shaded}

\begin{verbatim}
[1] 8
[1] 6
[1] 4
\end{verbatim}

\begin{Shaded}
\begin{Highlighting}[]
\CommentTok{\# Example: Use break inside a for loop when condition met}
\ControlFlowTok{for}\NormalTok{ (i }\ControlFlowTok{in} \DecValTok{1}\SpecialCharTok{:}\DecValTok{10}\NormalTok{) \{}
  \ControlFlowTok{if}\NormalTok{ (i }\SpecialCharTok{==} \DecValTok{6}\NormalTok{) \{}
    \ControlFlowTok{break}  \CommentTok{\# exit loop when i is 6}
\NormalTok{  \}}
  \FunctionTok{print}\NormalTok{(i)}
\NormalTok{\}}
\end{Highlighting}
\end{Shaded}

\begin{verbatim}
[1] 1
[1] 2
[1] 3
[1] 4
[1] 5
\end{verbatim}

\begin{Shaded}
\begin{Highlighting}[]
\CommentTok{\# Example: Skip printing the number 3 in a loop}
\ControlFlowTok{for}\NormalTok{ (i }\ControlFlowTok{in} \DecValTok{1}\SpecialCharTok{:}\DecValTok{5}\NormalTok{) \{}
  \ControlFlowTok{if}\NormalTok{ (i }\SpecialCharTok{==} \DecValTok{3}\NormalTok{) \{}
    \ControlFlowTok{next}  \CommentTok{\# skip the rest of this loop cycle}
\NormalTok{  \}}
  \FunctionTok{print}\NormalTok{(i)}
\NormalTok{\}}
\end{Highlighting}
\end{Shaded}

\begin{verbatim}
[1] 1
[1] 2
[1] 4
[1] 5
\end{verbatim}

\begin{Shaded}
\begin{Highlighting}[]
\CommentTok{\#   A simple function that adds two numbers}
\NormalTok{add\_two }\OtherTok{\textless{}{-}} \ControlFlowTok{function}\NormalTok{(a, b) \{}
  \FunctionTok{return}\NormalTok{(a }\SpecialCharTok{+}\NormalTok{ b)}
\NormalTok{\}}
\FunctionTok{print}\NormalTok{(}\FunctionTok{add\_two}\NormalTok{(}\DecValTok{3}\NormalTok{, }\DecValTok{5}\NormalTok{))}
\end{Highlighting}
\end{Shaded}

\begin{verbatim}
[1] 8
\end{verbatim}

\begin{Shaded}
\begin{Highlighting}[]
\CommentTok{\#   Function with a default parameter}
\NormalTok{greet }\OtherTok{\textless{}{-}} \ControlFlowTok{function}\NormalTok{(}\AttributeTok{name =} \StringTok{"friend"}\NormalTok{) \{}
  \FunctionTok{paste}\NormalTok{(}\StringTok{"Hello"}\NormalTok{, name)}
\NormalTok{\}}
\FunctionTok{print}\NormalTok{(}\FunctionTok{greet}\NormalTok{(}\StringTok{"Alice"}\NormalTok{))}
\end{Highlighting}
\end{Shaded}

\begin{verbatim}
[1] "Hello Alice"
\end{verbatim}

\begin{Shaded}
\begin{Highlighting}[]
\FunctionTok{print}\NormalTok{(}\FunctionTok{greet}\NormalTok{())}
\end{Highlighting}
\end{Shaded}

\begin{verbatim}
[1] "Hello friend"
\end{verbatim}

\begin{Shaded}
\begin{Highlighting}[]
\CommentTok{\#   A simple function that adds two numbers}
\NormalTok{add\_two }\OtherTok{\textless{}{-}} \ControlFlowTok{function}\NormalTok{(a, b) \{}
  \FunctionTok{return}\NormalTok{(a }\SpecialCharTok{+}\NormalTok{ b)}
\NormalTok{\}}
\FunctionTok{print}\NormalTok{(}\FunctionTok{add\_two}\NormalTok{(}\DecValTok{3}\NormalTok{, }\DecValTok{5}\NormalTok{))}
\end{Highlighting}
\end{Shaded}

\begin{verbatim}
[1] 8
\end{verbatim}

\begin{Shaded}
\begin{Highlighting}[]
\CommentTok{\#   Function with a default parameter}
\NormalTok{greet }\OtherTok{\textless{}{-}} \ControlFlowTok{function}\NormalTok{(}\AttributeTok{name =} \StringTok{"friend"}\NormalTok{) \{}
  \FunctionTok{paste}\NormalTok{(}\StringTok{"Hello"}\NormalTok{, name)}
\NormalTok{\}}
\FunctionTok{print}\NormalTok{(}\FunctionTok{greet}\NormalTok{(}\StringTok{"Alice"}\NormalTok{))}
\end{Highlighting}
\end{Shaded}

\begin{verbatim}
[1] "Hello Alice"
\end{verbatim}

\begin{Shaded}
\begin{Highlighting}[]
\FunctionTok{print}\NormalTok{(}\FunctionTok{greet}\NormalTok{())}
\end{Highlighting}
\end{Shaded}

\begin{verbatim}
[1] "Hello friend"
\end{verbatim}

\begin{Shaded}
\begin{Highlighting}[]
\CommentTok{\# (Used above in the add\_two and greet functions)}
\end{Highlighting}
\end{Shaded}

\begin{Shaded}
\begin{Highlighting}[]
\CommentTok{\#   Stop if an invalid condition is met inside a function}
\NormalTok{check\_positive }\OtherTok{\textless{}{-}} \ControlFlowTok{function}\NormalTok{(x) \{}
  \ControlFlowTok{if}\NormalTok{ (x }\SpecialCharTok{\textless{}} \DecValTok{0}\NormalTok{) \{}
    \FunctionTok{stop}\NormalTok{(}\StringTok{"x must be non{-}negative!"}\NormalTok{)}
\NormalTok{  \}}
  \FunctionTok{return}\NormalTok{(x)}
\NormalTok{\}}
\CommentTok{\# Uncomment the following line to see the error:}
\CommentTok{\# check\_positive({-}5)}

\CommentTok{\#   Stop execution manually}
\CommentTok{\# Uncomment the next line to stop the script:}
\CommentTok{\# stop("Manual stop here!")}
\end{Highlighting}
\end{Shaded}

\begin{Shaded}
\begin{Highlighting}[]
\CommentTok{\#   Issue a warning if a value is too high}
\NormalTok{warn\_if\_high }\OtherTok{\textless{}{-}} \ControlFlowTok{function}\NormalTok{(x) \{}
  \ControlFlowTok{if}\NormalTok{ (x }\SpecialCharTok{\textgreater{}} \DecValTok{100}\NormalTok{) \{}
    \FunctionTok{warning}\NormalTok{(}\StringTok{"x is very high!"}\NormalTok{)}
\NormalTok{  \}}
  \FunctionTok{return}\NormalTok{(x)}
\NormalTok{\}}
\FunctionTok{print}\NormalTok{(}\FunctionTok{warn\_if\_high}\NormalTok{(}\DecValTok{150}\NormalTok{))}
\end{Highlighting}
\end{Shaded}

\begin{verbatim}
Warning in warn_if_high(150): x is very high!
\end{verbatim}

\begin{verbatim}
[1] 150
\end{verbatim}

\begin{Shaded}
\begin{Highlighting}[]
\FunctionTok{print}\NormalTok{(}\FunctionTok{warn\_if\_high}\NormalTok{(}\DecValTok{50}\NormalTok{))}
\end{Highlighting}
\end{Shaded}

\begin{verbatim}
[1] 50
\end{verbatim}

\begin{Shaded}
\begin{Highlighting}[]
\CommentTok{\#   Try to take the square root of a negative number}
\NormalTok{result\_try1 }\OtherTok{\textless{}{-}} \FunctionTok{try}\NormalTok{(}\FunctionTok{sqrt}\NormalTok{(}\SpecialCharTok{{-}}\DecValTok{1}\NormalTok{))}
\end{Highlighting}
\end{Shaded}

\begin{verbatim}
Warning in sqrt(-1): NaNs produced
\end{verbatim}

\begin{Shaded}
\begin{Highlighting}[]
\FunctionTok{print}\NormalTok{(result\_try1)}
\end{Highlighting}
\end{Shaded}

\begin{verbatim}
[1] NaN
\end{verbatim}

\begin{Shaded}
\begin{Highlighting}[]
\CommentTok{\#   Try reading a non{-}existent variable}
\NormalTok{result\_try2 }\OtherTok{\textless{}{-}} \FunctionTok{try}\NormalTok{(}\FunctionTok{print}\NormalTok{(nonexistent\_variable))}
\end{Highlighting}
\end{Shaded}

\begin{verbatim}
Error in eval(expr, envir) : object 'nonexistent_variable' not found
\end{verbatim}

\begin{Shaded}
\begin{Highlighting}[]
\FunctionTok{print}\NormalTok{(result\_try2)}
\end{Highlighting}
\end{Shaded}

\begin{verbatim}
[1] "Error in eval(expr, envir) : object 'nonexistent_variable' not found\n"
attr(,"class")
[1] "try-error"
attr(,"condition")
<simpleError in eval(expr, envir): object 'nonexistent_variable' not found>
\end{verbatim}

\begin{Shaded}
\begin{Highlighting}[]
\CommentTok{\#   Use tryCatch to handle an error}
\NormalTok{result\_tc }\OtherTok{\textless{}{-}} \FunctionTok{tryCatch}\NormalTok{(\{}
  \FunctionTok{log}\NormalTok{(}\StringTok{"a"}\NormalTok{)  }\CommentTok{\# will cause an error}
\NormalTok{\}, }\AttributeTok{error =} \ControlFlowTok{function}\NormalTok{(e) \{}
  \StringTok{"An error occurred"}
\NormalTok{\})}
\FunctionTok{print}\NormalTok{(result\_tc)}
\end{Highlighting}
\end{Shaded}

\begin{verbatim}
[1] "An error occurred"
\end{verbatim}

\begin{Shaded}
\begin{Highlighting}[]
\CommentTok{\#   tryCatch with finally block}
\NormalTok{result\_tc2 }\OtherTok{\textless{}{-}} \FunctionTok{tryCatch}\NormalTok{(\{}
  \FunctionTok{stop}\NormalTok{(}\StringTok{"Oops!"}\NormalTok{)}
\NormalTok{\}, }\AttributeTok{error =} \ControlFlowTok{function}\NormalTok{(e) \{}
  \StringTok{"Caught the error"}
\NormalTok{\}, }\AttributeTok{finally =}\NormalTok{ \{}
  \FunctionTok{cat}\NormalTok{(}\StringTok{"This is the cleanup code.}\SpecialCharTok{\textbackslash{}n}\StringTok{"}\NormalTok{)}
\NormalTok{\})}
\end{Highlighting}
\end{Shaded}

\begin{verbatim}
This is the cleanup code.
\end{verbatim}

\begin{Shaded}
\begin{Highlighting}[]
\FunctionTok{print}\NormalTok{(result\_tc2)}
\end{Highlighting}
\end{Shaded}

\begin{verbatim}
[1] "Caught the error"
\end{verbatim}

\section{Working with Files}\label{5-working-with-files}

\paragraph{getwd() -- Get current working
directory}\label{getwd-get-current-working-directory}

\begin{Shaded}
\begin{Highlighting}[]
\CommentTok{\#   Simply print the working directory}
\NormalTok{wd1 }\OtherTok{\textless{}{-}} \FunctionTok{getwd}\NormalTok{()}
\FunctionTok{print}\NormalTok{(wd1)}
\end{Highlighting}
\end{Shaded}

\begin{verbatim}
[1] "D:/Sem-2/R-/3"
\end{verbatim}

\begin{Shaded}
\begin{Highlighting}[]
\CommentTok{\#   Save and print working directory from a variable}
\NormalTok{wd2 }\OtherTok{\textless{}{-}} \FunctionTok{getwd}\NormalTok{()}
\FunctionTok{cat}\NormalTok{(}\StringTok{"Working directory is:"}\NormalTok{, wd2, }\StringTok{"}\SpecialCharTok{\textbackslash{}n}\StringTok{"}\NormalTok{)}
\end{Highlighting}
\end{Shaded}

\begin{verbatim}
Working directory is: D:/Sem-2/R-/3 
\end{verbatim}

\paragraph{setwd() -- Set the working
directory}\label{setwd-set-the-working-directory}

\begin{Shaded}
\begin{Highlighting}[]
\CommentTok{\# (Be careful with this! It changes your folder for reading/writing files.)}
\CommentTok{\#   Change to a temporary directory (this should work safely)}
\CommentTok{\# setwd(tempdir())}
\CommentTok{\# cat("Now in temporary directory:", getwd(), "\textbackslash{}n")}

\CommentTok{\#   Set back to original (if you saved it before)}
\CommentTok{\# original\_wd \textless{}{-} "your/original/path"  \# replace with your actual path}
\CommentTok{\# setwd(original\_wd)}
\end{Highlighting}
\end{Shaded}

\paragraph{list.files() -- List files in a
directory}\label{list.files-list-files-in-a-directory}

\begin{Shaded}
\begin{Highlighting}[]
\CommentTok{\#   List files in current directory}
\NormalTok{files1 }\OtherTok{\textless{}{-}} \FunctionTok{list.files}\NormalTok{()}
\FunctionTok{print}\NormalTok{(files1)}
\end{Highlighting}
\end{Shaded}

\begin{verbatim}
[1] "Assignment-3.pdf"      "assignment3.qmd"       "assignment3.rmarkdown"
[4] "sample_data.rds"       "sample_output.csv"     "sample_output2.csv"   
[7] "sample_table.txt"      "sample_table2.txt"     "sample_vector.rds"    
\end{verbatim}

\begin{Shaded}
\begin{Highlighting}[]
\CommentTok{\#   List files with a specific pattern (e.g., files ending with .R)}
\NormalTok{files2 }\OtherTok{\textless{}{-}} \FunctionTok{list.files}\NormalTok{(}\AttributeTok{pattern =} \StringTok{"}\SpecialCharTok{\textbackslash{}\textbackslash{}}\StringTok{.R$"}\NormalTok{)}
\FunctionTok{print}\NormalTok{(files2)}
\end{Highlighting}
\end{Shaded}

\begin{verbatim}
character(0)
\end{verbatim}

\paragraph{file.exists() -- Check if a file
exists}\label{fileexists--check-if-a-file-exists}

\begin{Shaded}
\begin{Highlighting}[]
\CommentTok{\#   Check for a file named "test.txt"}
\NormalTok{exists1 }\OtherTok{\textless{}{-}} \FunctionTok{file.exists}\NormalTok{(}\StringTok{"test.txt"}\NormalTok{)}
\FunctionTok{print}\NormalTok{(exists1)}
\end{Highlighting}
\end{Shaded}

\begin{verbatim}
[1] FALSE
\end{verbatim}

\begin{Shaded}
\begin{Highlighting}[]
\CommentTok{\#   Check multiple files at once}
\NormalTok{exists2 }\OtherTok{\textless{}{-}} \FunctionTok{file.exists}\NormalTok{(}\FunctionTok{c}\NormalTok{(}\StringTok{"test.txt"}\NormalTok{, }\StringTok{"nonexistent.txt"}\NormalTok{))}
\FunctionTok{print}\NormalTok{(exists2)}
\end{Highlighting}
\end{Shaded}

\begin{verbatim}
[1] FALSE FALSE
\end{verbatim}

\paragraph{file.remove() -- Delete a
file}\label{file.remove-delete-a-file}

\begin{Shaded}
\begin{Highlighting}[]
\CommentTok{\#   Create and then remove a temporary file}
\NormalTok{temp\_file }\OtherTok{\textless{}{-}} \FunctionTok{tempfile}\NormalTok{(}\AttributeTok{fileext =} \StringTok{".txt"}\NormalTok{)}
\FunctionTok{cat}\NormalTok{(}\StringTok{"Temporary content"}\NormalTok{, }\AttributeTok{file =}\NormalTok{ temp\_file)}
\NormalTok{removed1 }\OtherTok{\textless{}{-}} \FunctionTok{file.remove}\NormalTok{(temp\_file)}
\FunctionTok{print}\NormalTok{(removed1)}
\end{Highlighting}
\end{Shaded}

\begin{verbatim}
[1] TRUE
\end{verbatim}

\begin{Shaded}
\begin{Highlighting}[]
\CommentTok{\#   Try removing a file that likely does not exist}
\NormalTok{removed2 }\OtherTok{\textless{}{-}} \FunctionTok{file.remove}\NormalTok{(}\StringTok{"no\_file.txt"}\NormalTok{)}
\end{Highlighting}
\end{Shaded}

\begin{verbatim}
Warning in file.remove("no_file.txt"): cannot remove file 'no_file.txt', reason
'No such file or directory'
\end{verbatim}

\begin{Shaded}
\begin{Highlighting}[]
\FunctionTok{print}\NormalTok{(removed2)}
\end{Highlighting}
\end{Shaded}

\begin{verbatim}
[1] FALSE
\end{verbatim}

\paragraph{read.csv() -- Read a CSV
file}\label{read.csv-read-a-csv-file}

\begin{Shaded}
\begin{Highlighting}[]
\CommentTok{\#   Write a small CSV and read it back}
\NormalTok{df\_sample }\OtherTok{\textless{}{-}} \FunctionTok{data.frame}\NormalTok{(}\AttributeTok{Name =} \FunctionTok{c}\NormalTok{(}\StringTok{"A"}\NormalTok{, }\StringTok{"B"}\NormalTok{), }\AttributeTok{Value =} \FunctionTok{c}\NormalTok{(}\DecValTok{10}\NormalTok{, }\DecValTok{20}\NormalTok{))}
\NormalTok{temp\_csv }\OtherTok{\textless{}{-}} \FunctionTok{tempfile}\NormalTok{(}\AttributeTok{fileext =} \StringTok{".csv"}\NormalTok{)}
\FunctionTok{write.csv}\NormalTok{(df\_sample, temp\_csv, }\AttributeTok{row.names =} \ConstantTok{FALSE}\NormalTok{)}
\NormalTok{df\_read1 }\OtherTok{\textless{}{-}} \FunctionTok{read.csv}\NormalTok{(temp\_csv)}
\FunctionTok{print}\NormalTok{(df\_read1)}
\end{Highlighting}
\end{Shaded}

\begin{verbatim}
  Name Value
1    A    10
2    B    20
\end{verbatim}

\begin{Shaded}
\begin{Highlighting}[]
\CommentTok{\#   Read CSV with header (default behavior)}
\NormalTok{df\_read2 }\OtherTok{\textless{}{-}} \FunctionTok{read.csv}\NormalTok{(temp\_csv, }\AttributeTok{header =} \ConstantTok{TRUE}\NormalTok{)}
\FunctionTok{print}\NormalTok{(df\_read2)}
\end{Highlighting}
\end{Shaded}

\begin{verbatim}
  Name Value
1    A    10
2    B    20
\end{verbatim}

\paragraph{write.csv() -- Write data to a CSV
file}\label{write.csv-write-data-to-a-csv-file}

\begin{Shaded}
\begin{Highlighting}[]
\CommentTok{\#   Write a data frame to CSV}
\FunctionTok{write.csv}\NormalTok{(df\_sample, }\StringTok{"sample\_output.csv"}\NormalTok{, }\AttributeTok{row.names =} \ConstantTok{TRUE}\NormalTok{)}
\CommentTok{\# Check if file exists afterward}
\FunctionTok{print}\NormalTok{(}\FunctionTok{file.exists}\NormalTok{(}\StringTok{"sample\_output.csv"}\NormalTok{))}
\end{Highlighting}
\end{Shaded}

\begin{verbatim}
[1] TRUE
\end{verbatim}

\begin{Shaded}
\begin{Highlighting}[]
\CommentTok{\#   Write CSV without row names}
\FunctionTok{write.csv}\NormalTok{(df\_sample, }\StringTok{"sample\_output2.csv"}\NormalTok{, }\AttributeTok{row.names =} \ConstantTok{FALSE}\NormalTok{)}
\FunctionTok{print}\NormalTok{(}\FunctionTok{file.exists}\NormalTok{(}\StringTok{"sample\_output2.csv"}\NormalTok{))}
\end{Highlighting}
\end{Shaded}

\begin{verbatim}
[1] TRUE
\end{verbatim}

\paragraph{read.table() -- Read tabular data from a
file}\label{read.table-read-tabular-data-from-a-file}

\begin{Shaded}
\begin{Highlighting}[]
\CommentTok{\#   Create a text file and read it as table}
\NormalTok{temp\_txt }\OtherTok{\textless{}{-}} \FunctionTok{tempfile}\NormalTok{(}\AttributeTok{fileext =} \StringTok{".txt"}\NormalTok{)}
\FunctionTok{write}\NormalTok{(}\StringTok{"1 2 3}\SpecialCharTok{\textbackslash{}n}\StringTok{4 5 6"}\NormalTok{, }\AttributeTok{file =}\NormalTok{ temp\_txt)}
\NormalTok{table1 }\OtherTok{\textless{}{-}} \FunctionTok{read.table}\NormalTok{(temp\_txt)}
\FunctionTok{print}\NormalTok{(table1)}
\end{Highlighting}
\end{Shaded}

\begin{verbatim}
  V1 V2 V3
1  1  2  3
2  4  5  6
\end{verbatim}

\begin{Shaded}
\begin{Highlighting}[]
\CommentTok{\#   Read table with a specified separator (space here)}
\NormalTok{table2 }\OtherTok{\textless{}{-}} \FunctionTok{read.table}\NormalTok{(temp\_txt, }\AttributeTok{sep =} \StringTok{" "}\NormalTok{)}
\FunctionTok{print}\NormalTok{(table2)}
\end{Highlighting}
\end{Shaded}

\begin{verbatim}
  V1 V2 V3
1  1  2  3
2  4  5  6
\end{verbatim}

\paragraph{write.table() -- Write data to a file in table
format}\label{write.table-write-data-to-a-file-in-table-format}

\begin{Shaded}
\begin{Highlighting}[]
\CommentTok{\#   Write a data frame to a text file}
\FunctionTok{write.table}\NormalTok{(df\_sample, }\StringTok{"sample\_table.txt"}\NormalTok{)}
\FunctionTok{print}\NormalTok{(}\FunctionTok{file.exists}\NormalTok{(}\StringTok{"sample\_table.txt"}\NormalTok{))}
\end{Highlighting}
\end{Shaded}

\begin{verbatim}
[1] TRUE
\end{verbatim}

\begin{Shaded}
\begin{Highlighting}[]
\CommentTok{\#   Write table with row names disabled}
\FunctionTok{write.table}\NormalTok{(df\_sample, }\StringTok{"sample\_table2.txt"}\NormalTok{, }\AttributeTok{row.names =} \ConstantTok{FALSE}\NormalTok{)}
\FunctionTok{print}\NormalTok{(}\FunctionTok{file.exists}\NormalTok{(}\StringTok{"sample\_table2.txt"}\NormalTok{))}
\end{Highlighting}
\end{Shaded}

\begin{verbatim}
[1] TRUE
\end{verbatim}

\paragraph{saveRDS() -- Save a single R object to a
file}\label{saverds-save-a-single-r-object-to-a-file}

\begin{Shaded}
\begin{Highlighting}[]
\CommentTok{\#   Save a data frame as an RDS file}
\FunctionTok{saveRDS}\NormalTok{(df\_sample, }\StringTok{"sample\_data.rds"}\NormalTok{)}
\FunctionTok{print}\NormalTok{(}\FunctionTok{file.exists}\NormalTok{(}\StringTok{"sample\_data.rds"}\NormalTok{))}
\end{Highlighting}
\end{Shaded}

\begin{verbatim}
[1] TRUE
\end{verbatim}

\begin{Shaded}
\begin{Highlighting}[]
\CommentTok{\#   Save a vector as an RDS file}
\FunctionTok{saveRDS}\NormalTok{(}\FunctionTok{c}\NormalTok{(}\DecValTok{1}\NormalTok{,}\DecValTok{2}\NormalTok{,}\DecValTok{3}\NormalTok{), }\StringTok{"sample\_vector.rds"}\NormalTok{)}
\FunctionTok{print}\NormalTok{(}\FunctionTok{file.exists}\NormalTok{(}\StringTok{"sample\_vector.rds"}\NormalTok{))}
\end{Highlighting}
\end{Shaded}

\begin{verbatim}
[1] TRUE
\end{verbatim}

\paragraph{readRDS() -- Read a single R object from an RDS
file}\label{readrds-read-a-single-r-object-from-an-rds-file}

\begin{Shaded}
\begin{Highlighting}[]
\CommentTok{\#   Read the previously saved data frame}
\NormalTok{df\_from\_rds }\OtherTok{\textless{}{-}} \FunctionTok{readRDS}\NormalTok{(}\StringTok{"sample\_data.rds"}\NormalTok{)}
\FunctionTok{print}\NormalTok{(df\_from\_rds)}
\end{Highlighting}
\end{Shaded}

\begin{verbatim}
  Name Value
1    A    10
2    B    20
\end{verbatim}

\begin{Shaded}
\begin{Highlighting}[]
\CommentTok{\#   Read the vector from the RDS file}
\NormalTok{vec\_from\_rds }\OtherTok{\textless{}{-}} \FunctionTok{readRDS}\NormalTok{(}\StringTok{"sample\_vector.rds"}\NormalTok{)}
\FunctionTok{print}\NormalTok{(vec\_from\_rds)}
\end{Highlighting}
\end{Shaded}

\begin{verbatim}
[1] 1 2 3
\end{verbatim}

\section{Working with Factors}\label{6-working-with-factors}

\paragraph{factor() -- Create a factor
variable}\label{factor-create-a-factor-variable}

\begin{Shaded}
\begin{Highlighting}[]
\CommentTok{\#   Convert a character vector to a factor}
\NormalTok{fact\_a }\OtherTok{\textless{}{-}} \FunctionTok{factor}\NormalTok{(}\FunctionTok{c}\NormalTok{(}\StringTok{"low"}\NormalTok{, }\StringTok{"medium"}\NormalTok{, }\StringTok{"high"}\NormalTok{, }\StringTok{"medium"}\NormalTok{))}
\FunctionTok{print}\NormalTok{(fact\_a)}
\end{Highlighting}
\end{Shaded}

\begin{verbatim}
[1] low    medium high   medium
Levels: high low medium
\end{verbatim}

\begin{Shaded}
\begin{Highlighting}[]
\CommentTok{\#   Specify levels manually}
\NormalTok{fact\_b }\OtherTok{\textless{}{-}} \FunctionTok{factor}\NormalTok{(}\FunctionTok{c}\NormalTok{(}\StringTok{"red"}\NormalTok{, }\StringTok{"blue"}\NormalTok{, }\StringTok{"red"}\NormalTok{), }\AttributeTok{levels =} \FunctionTok{c}\NormalTok{(}\StringTok{"red"}\NormalTok{, }\StringTok{"blue"}\NormalTok{, }\StringTok{"green"}\NormalTok{))}
\FunctionTok{print}\NormalTok{(fact\_b)}
\end{Highlighting}
\end{Shaded}

\begin{verbatim}
[1] red  blue red 
Levels: red blue green
\end{verbatim}

\paragraph{levels() -- Get or set levels of a
factor}\label{levels-get-or-set-levels-of-a-factor}

\begin{Shaded}
\begin{Highlighting}[]
\CommentTok{\#   Get levels of a factor}
\NormalTok{levels\_a }\OtherTok{\textless{}{-}} \FunctionTok{levels}\NormalTok{(fact\_a)}
\FunctionTok{print}\NormalTok{(levels\_a)}
\end{Highlighting}
\end{Shaded}

\begin{verbatim}
[1] "high"   "low"    "medium"
\end{verbatim}

\begin{Shaded}
\begin{Highlighting}[]
\CommentTok{\#   Change levels of a factor}
\FunctionTok{levels}\NormalTok{(fact\_a) }\OtherTok{\textless{}{-}} \FunctionTok{c}\NormalTok{(}\StringTok{"L"}\NormalTok{, }\StringTok{"M"}\NormalTok{, }\StringTok{"H"}\NormalTok{)}
\FunctionTok{print}\NormalTok{(fact\_a)}
\end{Highlighting}
\end{Shaded}

\begin{verbatim}
[1] M H L H
Levels: L M H
\end{verbatim}

\paragraph{relevel() -- Change the reference level of a
factor}\label{relevel-change-the-reference-level-of-a-factor}

\begin{Shaded}
\begin{Highlighting}[]
\CommentTok{\#   Make "blue" the reference level}
\NormalTok{fact\_c }\OtherTok{\textless{}{-}} \FunctionTok{factor}\NormalTok{(}\FunctionTok{c}\NormalTok{(}\StringTok{"red"}\NormalTok{, }\StringTok{"blue"}\NormalTok{, }\StringTok{"red"}\NormalTok{, }\StringTok{"green"}\NormalTok{))}
\NormalTok{fact\_c }\OtherTok{\textless{}{-}} \FunctionTok{relevel}\NormalTok{(fact\_c, }\AttributeTok{ref =} \StringTok{"blue"}\NormalTok{)}
\FunctionTok{print}\NormalTok{(fact\_c)}
\end{Highlighting}
\end{Shaded}

\begin{verbatim}
[1] red   blue  red   green
Levels: blue green red
\end{verbatim}

\begin{Shaded}
\begin{Highlighting}[]
\CommentTok{\#   Relevel with a different reference}
\NormalTok{fact\_d }\OtherTok{\textless{}{-}} \FunctionTok{factor}\NormalTok{(}\FunctionTok{c}\NormalTok{(}\StringTok{"A"}\NormalTok{, }\StringTok{"B"}\NormalTok{, }\StringTok{"C"}\NormalTok{, }\StringTok{"A"}\NormalTok{))}
\NormalTok{fact\_d }\OtherTok{\textless{}{-}} \FunctionTok{relevel}\NormalTok{(fact\_d, }\AttributeTok{ref =} \StringTok{"C"}\NormalTok{)}
\FunctionTok{print}\NormalTok{(fact\_d)}
\end{Highlighting}
\end{Shaded}

\begin{verbatim}
[1] A B C A
Levels: C A B
\end{verbatim}

\paragraph{cut() -- Divide numeric data into intervals
(factors)}\label{cut-divide-numeric-data-into-intervals-factors}

\begin{Shaded}
\begin{Highlighting}[]
\CommentTok{\#   Cut numbers into equal{-}width bins}
\NormalTok{nums }\OtherTok{\textless{}{-}} \FunctionTok{c}\NormalTok{(}\DecValTok{1}\NormalTok{, }\DecValTok{3}\NormalTok{, }\DecValTok{5}\NormalTok{, }\DecValTok{7}\NormalTok{, }\DecValTok{9}\NormalTok{)}
\NormalTok{cut1 }\OtherTok{\textless{}{-}} \FunctionTok{cut}\NormalTok{(nums, }\AttributeTok{breaks =} \DecValTok{3}\NormalTok{)}
\FunctionTok{print}\NormalTok{(cut1)}
\end{Highlighting}
\end{Shaded}

\begin{verbatim}
[1] (0.992,3.67] (0.992,3.67] (3.67,6.33]  (6.33,9.01]  (6.33,9.01] 
Levels: (0.992,3.67] (3.67,6.33] (6.33,9.01]
\end{verbatim}

\begin{Shaded}
\begin{Highlighting}[]
\CommentTok{\#   Cut with custom labels}
\NormalTok{cut2 }\OtherTok{\textless{}{-}} \FunctionTok{cut}\NormalTok{(nums, }\AttributeTok{breaks =} \FunctionTok{c}\NormalTok{(}\DecValTok{0}\NormalTok{, }\DecValTok{4}\NormalTok{, }\DecValTok{8}\NormalTok{, }\DecValTok{10}\NormalTok{), }\AttributeTok{labels =} \FunctionTok{c}\NormalTok{(}\StringTok{"low"}\NormalTok{, }\StringTok{"medium"}\NormalTok{, }\StringTok{"high"}\NormalTok{))}
\FunctionTok{print}\NormalTok{(cut2)}
\end{Highlighting}
\end{Shaded}

\begin{verbatim}
[1] low    low    medium medium high  
Levels: low medium high
\end{verbatim}

\section{Date \& Time Handling}\label{7-date--time-handling}

\paragraph{as.POSIXct() -- Convert string to
date-time}\label{as.posixct-convert-string-to-date-time}

\begin{Shaded}
\begin{Highlighting}[]
\CommentTok{\#   Convert a date{-}time string}
\NormalTok{dt1 }\OtherTok{\textless{}{-}} \FunctionTok{as.POSIXct}\NormalTok{(}\StringTok{"2025{-}02{-}03 12:34:56"}\NormalTok{)}
\FunctionTok{print}\NormalTok{(dt1)}
\end{Highlighting}
\end{Shaded}

\begin{verbatim}
[1] "2025-02-03 12:34:56 IST"
\end{verbatim}

\begin{Shaded}
\begin{Highlighting}[]
\CommentTok{\#   Convert with a specific format}
\NormalTok{dt2 }\OtherTok{\textless{}{-}} \FunctionTok{as.POSIXct}\NormalTok{(}\StringTok{"03/02/2025 12:34:56"}\NormalTok{, }\AttributeTok{format =} \StringTok{"\%d/\%m/\%Y \%H:\%M:\%S"}\NormalTok{)}
\FunctionTok{print}\NormalTok{(dt2)}
\end{Highlighting}
\end{Shaded}

\begin{verbatim}
[1] "2025-02-03 12:34:56 IST"
\end{verbatim}

\paragraph{as.POSIXlt() -- Convert string to list-like
date-time}\label{as.posixlt-convert-string-to-list-like-date-time}

\begin{Shaded}
\begin{Highlighting}[]
\CommentTok{\#   Convert and see the list components}
\NormalTok{dt\_lt1 }\OtherTok{\textless{}{-}} \FunctionTok{as.POSIXlt}\NormalTok{(}\StringTok{"2025{-}02{-}03 12:34:56"}\NormalTok{)}
\FunctionTok{print}\NormalTok{(dt\_lt1)}
\end{Highlighting}
\end{Shaded}

\begin{verbatim}
[1] "2025-02-03 12:34:56 IST"
\end{verbatim}

\begin{Shaded}
\begin{Highlighting}[]
\CommentTok{\#   Extract the year from the POSIXlt object}
\NormalTok{year\_val }\OtherTok{\textless{}{-}}\NormalTok{ dt\_lt1}\SpecialCharTok{$}\NormalTok{year }\SpecialCharTok{+} \DecValTok{1900}  \CommentTok{\# year is stored as years since 1900}
\FunctionTok{print}\NormalTok{(year\_val)}
\end{Highlighting}
\end{Shaded}

\begin{verbatim}
[1] 2025
\end{verbatim}

\paragraph{difftime() -- Compute the difference between two date-time
values}\label{difftime-compute-the-difference-between-two-date-time-values}

\begin{Shaded}
\begin{Highlighting}[]
\CommentTok{\#   Difference between two dates}
\NormalTok{diff1 }\OtherTok{\textless{}{-}} \FunctionTok{difftime}\NormalTok{(}\FunctionTok{as.Date}\NormalTok{(}\StringTok{"2025{-}02{-}10"}\NormalTok{), }\FunctionTok{as.Date}\NormalTok{(}\StringTok{"2025{-}02{-}03"}\NormalTok{))}
\FunctionTok{print}\NormalTok{(diff1)}
\end{Highlighting}
\end{Shaded}

\begin{verbatim}
Time difference of 7 days
\end{verbatim}

\begin{Shaded}
\begin{Highlighting}[]
\CommentTok{\#   Difference in hours between two date{-}times}
\NormalTok{dt\_start }\OtherTok{\textless{}{-}} \FunctionTok{as.POSIXct}\NormalTok{(}\StringTok{"2025{-}02{-}03 08:00:00"}\NormalTok{)}
\NormalTok{dt\_end   }\OtherTok{\textless{}{-}} \FunctionTok{as.POSIXct}\NormalTok{(}\StringTok{"2025{-}02{-}03 20:00:00"}\NormalTok{)}
\NormalTok{diff2 }\OtherTok{\textless{}{-}} \FunctionTok{difftime}\NormalTok{(dt\_end, dt\_start, }\AttributeTok{units =} \StringTok{"hours"}\NormalTok{)}
\FunctionTok{print}\NormalTok{(diff2)}
\end{Highlighting}
\end{Shaded}

\begin{verbatim}
Time difference of 12 hours
\end{verbatim}

\paragraph{format() -- Format date-time objects as
strings}\label{format-format-date-time-objects-as-strings}

\begin{Shaded}
\begin{Highlighting}[]
\CommentTok{\#   Format a date in a simple way}
\NormalTok{fmt1 }\OtherTok{\textless{}{-}} \FunctionTok{format}\NormalTok{(}\FunctionTok{as.Date}\NormalTok{(}\StringTok{"2025{-}02{-}03"}\NormalTok{), }\StringTok{"\%B \%d, \%Y"}\NormalTok{)}
\FunctionTok{print}\NormalTok{(fmt1)}
\end{Highlighting}
\end{Shaded}

\begin{verbatim}
[1] "February 03, 2025"
\end{verbatim}

\begin{Shaded}
\begin{Highlighting}[]
\CommentTok{\#   Format a date{-}time to show time and date}
\NormalTok{fmt2 }\OtherTok{\textless{}{-}} \FunctionTok{format}\NormalTok{(}\FunctionTok{Sys.Date}\NormalTok{(), }\StringTok{"\%Y/\%m/\%d"}\NormalTok{)}
\FunctionTok{print}\NormalTok{(fmt2)}
\end{Highlighting}
\end{Shaded}

\begin{verbatim}
[1] "2026/02/05"
\end{verbatim}

\section{Debugging \& Environment
Management}\label{debugging-environment-management}

\paragraph{ls() -- List objects in the
environment}\label{ls-list-objects-in-the-environment}

\begin{Shaded}
\begin{Highlighting}[]
\CommentTok{\#   List all current objects}
\NormalTok{objs1 }\OtherTok{\textless{}{-}} \FunctionTok{ls}\NormalTok{()}
\FunctionTok{print}\NormalTok{(objs1)}
\end{Highlighting}
\end{Shaded}

\begin{verbatim}
  [1] "a1"             "a2"             "a3"             "add_two"       
  [5] "al1"            "al2"            "al3"            "check_number"  
  [9] "check_positive" "count"          "counter"        "cut1"          
 [13] "cut2"           "df_from_rds"    "df_read1"       "df_read2"      
 [17] "df_sample"      "df1"            "df2"            "df3"           
 [21] "df4"            "df5"            "diff1"          "diff2"         
 [25] "dt_end"         "dt_lt1"         "dt_start"       "dt1"           
 [29] "dt2"            "exists1"        "exists2"        "fact_a"        
 [33] "fact_b"         "fact_c"         "fact_d"         "files1"        
 [37] "files2"         "fmt1"           "fmt2"           "greet"         
 [41] "i"              "in1"            "in2"            "in3"           
 [45] "l1"             "l2"             "l3"             "letter"        
 [49] "levels_a"       "lo1"            "lo2"            "lo3"           
 [53] "m1"             "m2"             "m3"             "mn1"           
 [57] "mn2"            "mn3"            "mnm1"           "mnm2"          
 [61] "mnm3"           "mx1"            "mx2"            "mx3"           
 [65] "na1"            "na2"            "na3"            "no1"           
 [69] "no2"            "no3"            "nums"           "o1"            
 [73] "o2"             "o3"             "p01"            "p02"           
 [77] "p03"            "p1"             "p2"             "p3"            
 [81] "r1"             "r2"             "r3"             "removed1"      
 [85] "removed2"       "rep1"           "rep2"           "rep3"          
 [89] "result_tc"      "result_tc2"     "result_try1"    "result_try2"   
 [93] "result1"        "result2"        "rp1"            "rp2"           
 [97] "rp3"            "s1"             "s2"             "s3"            
[101] "sb1"            "sb2"            "sb3"            "sm1"           
[105] "sm2"            "sm3"            "so1"            "so2"           
[109] "so3"            "t1"             "t2"             "t3"            
[113] "table1"         "table2"         "temp_csv"       "temp_file"     
[117] "temp_txt"       "u1"             "u2"             "u3"            
[121] "up1"            "up2"            "up3"            "v1"            
[125] "v2"             "v3"             "val"            "vec_from_rds"  
[129] "vec_loop"       "vec1"           "w1"             "w2"            
[133] "w3"             "warn_if_high"   "wd1"            "wd2"           
[137] "x"              "year_val"      
\end{verbatim}

\begin{Shaded}
\begin{Highlighting}[]
\CommentTok{\#   List objects that start with "df"}
\NormalTok{objs2 }\OtherTok{\textless{}{-}} \FunctionTok{ls}\NormalTok{(}\AttributeTok{pattern =} \StringTok{"\^{}df"}\NormalTok{)}
\FunctionTok{print}\NormalTok{(objs2)}
\end{Highlighting}
\end{Shaded}

\begin{verbatim}
[1] "df_from_rds" "df_read1"    "df_read2"    "df_sample"   "df1"        
[6] "df2"         "df3"         "df4"         "df5"        
\end{verbatim}

\paragraph{rm() -- Remove objects from the
environment}\label{rm-remove-objects-from-the-environment}

\begin{Shaded}
\begin{Highlighting}[]
\CommentTok{\#   Remove a single variable}
\NormalTok{temp\_var }\OtherTok{\textless{}{-}} \DecValTok{123}
\FunctionTok{print}\NormalTok{(temp\_var)}
\end{Highlighting}
\end{Shaded}

\begin{verbatim}
[1] 123
\end{verbatim}

\begin{Shaded}
\begin{Highlighting}[]
\FunctionTok{rm}\NormalTok{(temp\_var)}
\CommentTok{\# Uncommenting the next line would cause an error since temp\_var is removed:}
\CommentTok{\# print(temp\_var)}

\CommentTok{\#   Remove multiple objects}
\NormalTok{a }\OtherTok{\textless{}{-}} \DecValTok{1}\NormalTok{; b }\OtherTok{\textless{}{-}} \DecValTok{2}\NormalTok{; c }\OtherTok{\textless{}{-}} \DecValTok{3}
\FunctionTok{rm}\NormalTok{(a, b)}
\CommentTok{\# Only \textquotesingle{}c\textquotesingle{} should remain}
\FunctionTok{print}\NormalTok{(}\FunctionTok{ls}\NormalTok{())}
\end{Highlighting}
\end{Shaded}

\begin{verbatim}
  [1] "a1"             "a2"             "a3"             "add_two"       
  [5] "al1"            "al2"            "al3"            "c"             
  [9] "check_number"   "check_positive" "count"          "counter"       
 [13] "cut1"           "cut2"           "df_from_rds"    "df_read1"      
 [17] "df_read2"       "df_sample"      "df1"            "df2"           
 [21] "df3"            "df4"            "df5"            "diff1"         
 [25] "diff2"          "dt_end"         "dt_lt1"         "dt_start"      
 [29] "dt1"            "dt2"            "exists1"        "exists2"       
 [33] "fact_a"         "fact_b"         "fact_c"         "fact_d"        
 [37] "files1"         "files2"         "fmt1"           "fmt2"          
 [41] "greet"          "i"              "in1"            "in2"           
 [45] "in3"            "l1"             "l2"             "l3"            
 [49] "letter"         "levels_a"       "lo1"            "lo2"           
 [53] "lo3"            "m1"             "m2"             "m3"            
 [57] "mn1"            "mn2"            "mn3"            "mnm1"          
 [61] "mnm2"           "mnm3"           "mx1"            "mx2"           
 [65] "mx3"            "na1"            "na2"            "na3"           
 [69] "no1"            "no2"            "no3"            "nums"          
 [73] "o1"             "o2"             "o3"             "objs1"         
 [77] "objs2"          "p01"            "p02"            "p03"           
 [81] "p1"             "p2"             "p3"             "r1"            
 [85] "r2"             "r3"             "removed1"       "removed2"      
 [89] "rep1"           "rep2"           "rep3"           "result_tc"     
 [93] "result_tc2"     "result_try1"    "result_try2"    "result1"       
 [97] "result2"        "rp1"            "rp2"            "rp3"           
[101] "s1"             "s2"             "s3"             "sb1"           
[105] "sb2"            "sb3"            "sm1"            "sm2"           
[109] "sm3"            "so1"            "so2"            "so3"           
[113] "t1"             "t2"             "t3"             "table1"        
[117] "table2"         "temp_csv"       "temp_file"      "temp_txt"      
[121] "u1"             "u2"             "u3"             "up1"           
[125] "up2"            "up3"            "v1"             "v2"            
[129] "v3"             "val"            "vec_from_rds"   "vec_loop"      
[133] "vec1"           "w1"             "w2"             "w3"            
[137] "warn_if_high"   "wd1"            "wd2"            "x"             
[141] "year_val"      
\end{verbatim}

\paragraph{gc() -- Run garbage collection (free up
memory)}\label{gc-run-garbage-collection-free-up-memory}

\begin{Shaded}
\begin{Highlighting}[]
\CommentTok{\#   Run garbage collection and print result}
\NormalTok{gc\_result1 }\OtherTok{\textless{}{-}} \FunctionTok{gc}\NormalTok{()}
\FunctionTok{print}\NormalTok{(gc\_result1)}
\end{Highlighting}
\end{Shaded}

\begin{verbatim}
          used (Mb) gc trigger (Mb) max used (Mb)
Ncells  660910 35.3    1307455 69.9  1307455 69.9
Vcells 1245388  9.6    8388608 64.0  2568579 19.6
\end{verbatim}

\begin{Shaded}
\begin{Highlighting}[]
\CommentTok{\#   Just call gc() to see memory cleanup info}
\FunctionTok{gc}\NormalTok{()}
\end{Highlighting}
\end{Shaded}

\begin{verbatim}
          used (Mb) gc trigger (Mb) max used (Mb)
Ncells  660896 35.3    1307455 69.9  1307455 69.9
Vcells 1245380  9.6    8388608 64.0  2568579 19.6
\end{verbatim}

\paragraph{debug() -- Turn on debugging for a
function}\label{debug-turn-on-debugging-for-a-function}

\begin{Shaded}
\begin{Highlighting}[]
\CommentTok{\#   Define a simple function and set debug on it}
\NormalTok{test\_fun }\OtherTok{\textless{}{-}} \ControlFlowTok{function}\NormalTok{(x) \{}
\NormalTok{  y }\OtherTok{\textless{}{-}}\NormalTok{ x }\SpecialCharTok{+} \DecValTok{1}
  \FunctionTok{return}\NormalTok{(y)}
\NormalTok{\}}
\FunctionTok{debug}\NormalTok{(test\_fun)}
\CommentTok{\# Uncomment the next line to step through the function:}
\CommentTok{\# test\_fun(5)}
\FunctionTok{undebug}\NormalTok{(test\_fun)  }\CommentTok{\# Turn off debugging}
\end{Highlighting}
\end{Shaded}

\subsubsection{traceback() -- Show the call stack after an
error}\label{traceback-show-the-call-stack-after-an-error}

\begin{Shaded}
\begin{Highlighting}[]
\CommentTok{\#   Force an error then show traceback}
\NormalTok{error\_fun }\OtherTok{\textless{}{-}} \ControlFlowTok{function}\NormalTok{() \{}
  \FunctionTok{stop}\NormalTok{(}\StringTok{"This is an error!"}\NormalTok{)}
\NormalTok{\}}
\CommentTok{\# Uncomment the following lines to see traceback:}
\CommentTok{\# try(error\_fun())}
\CommentTok{\# traceback()}
\end{Highlighting}
\end{Shaded}

\paragraph{print() -- Print output to the
console}\label{print-print-output-to-the-console}

\begin{Shaded}
\begin{Highlighting}[]
\CommentTok{\#   Print a number}
\FunctionTok{print}\NormalTok{(}\DecValTok{100}\NormalTok{)}
\end{Highlighting}
\end{Shaded}

\begin{verbatim}
[1] 100
\end{verbatim}

\begin{Shaded}
\begin{Highlighting}[]
\CommentTok{\#   Print a message}
\FunctionTok{print}\NormalTok{(}\StringTok{"Hello, world!"}\NormalTok{)}
\end{Highlighting}
\end{Shaded}

\begin{verbatim}
[1] "Hello, world!"
\end{verbatim}

\begin{Shaded}
\begin{Highlighting}[]
\CommentTok{\# Extra example: Print a vector}
\FunctionTok{print}\NormalTok{(}\FunctionTok{c}\NormalTok{(}\StringTok{"a"}\NormalTok{, }\StringTok{"b"}\NormalTok{, }\StringTok{"c"}\NormalTok{))}
\end{Highlighting}
\end{Shaded}

\begin{verbatim}
[1] "a" "b" "c"
\end{verbatim}

\paragraph{cat() -- Concatenate and print text (no
quotes)}\label{cat-concatenate-and-print-text-no-quotes}

\begin{Shaded}
\begin{Highlighting}[]
\CommentTok{\#   Concatenate and print a simple message}
\FunctionTok{cat}\NormalTok{(}\StringTok{"This is a simple message.}\SpecialCharTok{\textbackslash{}n}\StringTok{"}\NormalTok{)}
\end{Highlighting}
\end{Shaded}

\begin{verbatim}
This is a simple message.
\end{verbatim}

\begin{Shaded}
\begin{Highlighting}[]
\CommentTok{\#   Concatenate multiple pieces of text}
\FunctionTok{cat}\NormalTok{(}\StringTok{"The answer is"}\NormalTok{, }\DecValTok{42}\NormalTok{, }\StringTok{"}\SpecialCharTok{\textbackslash{}n}\StringTok{"}\NormalTok{)}
\end{Highlighting}
\end{Shaded}

\begin{verbatim}
The answer is 42 
\end{verbatim}

\begin{Shaded}
\begin{Highlighting}[]
\CommentTok{\# Extra example: Print without automatic newlines (you add \textbackslash{}n manually)}
\FunctionTok{cat}\NormalTok{(}\StringTok{"Line one.}\SpecialCharTok{\textbackslash{}n}\StringTok{Line two.}\SpecialCharTok{\textbackslash{}n}\StringTok{"}\NormalTok{)}
\end{Highlighting}
\end{Shaded}

\begin{verbatim}
Line one.
Line two.
\end{verbatim}




\end{document}
